% \anonmode adds the acmart anonymous option and an anonymised title.
\ifdefined\anonmode
  \documentclass[acmsmall,anonymous]{acmart}
\else
  \documentclass[acmsmall]{acmart}
\fi

% acmart loads amsmath, amssymb, amsthm, microtype, xcolor, hyperref, natbib, enumitem.
\usepackage{array}
\usepackage{colortbl}
\DeclareMathAlphabet{\mathpzc}{OT1}{pzc}{m}{it}
\usepackage{mathpartir}
\usepackage{listings}
\usepackage{subcaption}
\usepackage{float}
\usepackage{xspace}
\usepackage{enumitem}
\usepackage{tikz}

% \finalmode renders the clean version, without change markup.
\ifdefined\finalmode\newcommand{\acceptchanges}{}\fi
\usepackage{changebar}
\ifdefined\finalmode\nochangebars\fi

% Document and references
\ifdefined\anonmode
  \newcommand*{\PurePy}{[Our lang]\xspace}
\else
  \newcommand*{\PurePy}{PurePy\xspace}
\fi
\ifdefined\anonmode
  \newcommand*{\issue}[1]{\##1}
\else
  \newcommand*{\issue}[1]{\href{https://github.com/pure-py/pure-py-spec/issues/#1}{\##1}}
\fi
\newcommand*{\versionstring}{}
\newcommand*{\version}[3]{%
  \gdef\versionstring{#1.#2.#3}%
}
\newcommand*{\figref}[1]{Figure~\ref{fig:#1}}
\newcommand*{\figrefTwo}[2]{Figures \ref{fig:#1} and \ref{fig:#2}}

% Figure and rule presentation
\definecolor{shadecolor}{rgb}{0.9,0.9,0.9}
\newcommand*{\shadebox}[1]{\fcolorbox{lightgray}{shadecolor}{\raisebox{0pt}[0.60\baselineskip][0.05\baselineskip]{#1}}}
\newcommand*{\ruleName}[1]{\textcolor{gray}{\textnormal{\textsf{#1}}}}

% General notation
\newcommand*{\set}[1]{\{#1\}}
\newcommand*{\dom}{\mathrm{dom}}
\newcommand*{\seq}[1]{\vec{#1}}
\newcommand*{\cons}{\cdot}
\newcommand*{\B}{\mathbb{B}}

% Terms: expressions and comprehension qualifiers
\newcommand*{\exCall}[2]{#1(#2)}
\newcommand*{\exConstr}[2]{#1(#2)}
\newcommand*{\exBinOp}[3]{#1 \mathbin{#2} #3}
\newcommand*{\exUnaryOp}[2]{#1\,#2}
\newcommand*{\exAttr}[2]{#1\mathord{.}#2}
\newcommand*{\exSubscript}[2]{#1[#2]}
\newcommand*{\exCondExpr}[3]{#1\;\pyinline{if}\;#2\;\pyinline{else}\;#3}
\newcommand*{\exList}[1]{[#1]}
\newcommand*{\exTuple}[1]{(#1)}
\newcommand*{\exDict}[1]{\{#1\}}
\newcommand*{\exListComp}[2]{[\,#1\;#2]}
\newcommand*{\exLambda}[2]{\pyinline{lambda}\;#1\mathord{:}\;#2}
\newcommand*{\qualGuard}[1]{\pyinline{if}\;#1}
\newcommand*{\qualGenerator}[2]{\pyinline{for}\;#1\;\pyinline{in}\;#2}

% Terms: statements, blocks, and module bodies
\newcommand*{\exPass}{\pyinline{pass}}
\newcommand*{\exAssign}[2]{#1 = #2}
\newcommand*{\exAssert}[1]{\pyinline{assert}\;#1}
\newcommand*{\exAssertMsg}[2]{\pyinline{assert}\;#1, #2}
\newcommand*{\exReturn}[1]{\pyinline{return}\;#1}
\newcommand*{\exReturnNone}{\pyinline{return}}
\newcommand*{\exIf}[3]{\pyinline{if}\;#1:\;#2\;#3}
\newcommand*{\exIfElse}[4]{\pyinline{if}\;#1:\;#2\;#3\;\pyinline{else}:\;#4}
\newcommand*{\exElif}[2]{\pyinline{elif}\;#1:\;#2}
\newcommand*{\exMatch}[2]{\pyinline{match}\;#1\mathord{:}\;#2}
\newcommand*{\exCase}[2]{\pyinline{case}\;#1\mathord{:}\;#2}
\newcommand*{\exDef}[3]{\pyinline{def}\;#1(#2):\;#3}
\newcommand*{\exDataclass}[2]{@\pyinline{dataclass}\;\pyinline{class}\;#1:\;#2}
\newcommand*{\exDataclassExtend}[3]{@\pyinline{dataclass}\;\pyinline{class}\;#1(#2):\;#3}
\newcommand*{\exField}[1]{#1:\;\pyinline{Any}}
\newcommand*{\exImport}[1]{\pyinline{import}\;#1}
\newcommand*{\exFromImport}[2]{\pyinline{from}\;#1\;\pyinline{import}\;#2}
\newcommand*{\exSeq}[2]{#1\;#2}

% Terms: patterns
\newcommand*{\patConstr}[2]{#1(#2)}
\newcommand*{\patWild}{\text{\pytext{\textunderscore}}}
\newcommand*{\patList}[1]{[#1]}
\newcommand*{\patTuple}[1]{(#1)}

% Values and context entries
\newcommand*{\bind}[2]{#1\mathord{:}\,#2}
\newcommand*{\statusDefAssigned}{\mathsf{tt}}
\newcommand*{\statusNotDefAssigned}{\mathsf{ff}}
\newcommand*{\modStub}[1]{\mathsf{Mod}(#1)}
\newcommand*{\modLoaded}[2]{\mathsf{Mod}(#1, #2)}
\newcommand*{\clsEntry}[4]{\mathsf{Cls}_{#1}(#2, #3, #4)}
\newcommand*{\exLamClosure}[4][]{\mathsf{Lam}_{#1}(#2, #3, #4)}
\newcommand*{\exDefClosure}[4][]{\mathsf{Def}_{#1}(#2, #3, #4)}
\newcommand*{\exObj}[3]{\mathsf{Obj}_{#1}(#2, #3)}

% Context and environment operations
\newcommand*{\override}{\cdot}
\newcommand*{\merge}{\oplus}
\newcommand*{\meet}{\wedge}
\newcommand*{\extend}{\lhd}
\newcommand*{\disjunion}{\uplus}

% Result types and results
\newcommand*{\Assigns}[1]{\mathsf{Assigns}\;#1}
\newcommand*{\Returns}{\mathsf{Returns}}
\newcommand{\assigns}[1]{\mathsf{assigns}\;#1}
\newcommand{\returns}[1]{\mathsf{returns}\;#1}

% Metafunctions
\newcommand*{\assignsS}{\mathrm{assigns}}
\newcommand*{\capturesS}{\mathrm{captures}}
\newcommand*{\capturesE}{\mathrm{captures}_{\mathsf{e}}}
\newcommand*{\fvS}{\mathrm{fv}}
\newcommand*{\fvE}{\mathrm{fv}_{\mathsf{e}}}
\newcommand*{\binds}{\mathrm{binds}}
\newcommand*{\ownFields}{\mathrm{own\text{-}fields}}
\newcommand*{\fields}{\mathrm{fields}}
\newcommand*{\ancestors}{\mathrm{ancestors}}
\newcommand*{\imports}{\mathrm{imports}}
\newcommand*{\deps}{\mathrm{deps}}
\newcommand*{\submodules}{\mathrm{submodules}}
\newcommand*{\nameOf}{\mathrm{name}}
\newcommand*{\unwrap}{\mathrm{unwrap}}
\newcommand*{\dispatch}{\mathrm{dispatch}}
\newcommand*{\env}{\mathrm{env}}
\newcommand*{\resolveClass}{\mathrm{resolve\text{-}class}}

% Orders, relations, and judgement symbols
\newcommand*{\prefixOf}{\preceq}
\newcommand*{\properPrefixOf}{\prec}
\newcommand*{\subsumed}{\sqsubseteq}
\newcommand*{\matches}{\rightsquigarrow}
\newcommand*{\notmatches}{\not\rightsquigarrow}
\newcommand*{\names}{\mathrel{\mathsf{names}}}
\newcommand*{\importsR}{\mathrel{\mathsf{imports}}}
\newcommand*{\loadsTo}{\mathrel{\mathsf{loads\mbox{-}to}}}

% Open term typing arrow
\newcommand{\carrow}[1]{\xrightarrow{\; #1\; }}

% Change markup (changes + changebar packages).
% No strikethrough: ulem's \sout is neither line-breakable nor metric-neutral.
\colorlet{deletedgray}{gray!70}
\setdeletedmarkup{\textcolor{deletedgray}{#1}}
% Capitalised variants add a margin changebar; use the lowercase forms inside
% figures or math, where changebar spans misbehave. \leavevmode commits a run-in
% heading in the default colour before the markup colour starts.
\newcommand{\Added}[1]{\leavevmode\protect\cbstart\added{#1}\protect\cbend\xspace}
\newcommand{\Deleted}[1]{\leavevmode\protect\cbstart\deleted{#1}\protect\cbend\xspace}
% The control space after #1 separates the new text from the deleted text shown
% after it (a plain space at the start of #2 is eaten by the changes package).
\newcommand{\Replaced}[2]{\leavevmode\protect\cbstart\replaced{#1\ }{#2}\protect\cbend\xspace}
% Deleted display material (diagrams, figures) that cannot live inside \Deleted:
% greyed out with a changebar in draft mode, discarded in final mode.
\newenvironment{DeletedBlock}{\par\cbstart\color{deletedgray}}{\cbend\par}
% Added display material (tables, figures) that cannot live inside \Added: alignment
% and rule material is fragile inside the changes package's argument.
\newenvironment{AddedBlock}{\par\cbstart\color{blue}}{\cbend\par}
% Final mode bypasses the changes package: its \IfIsEmpty test expands the
% replacement text, and any macro with #1 in its body then causes a runaway \iffalse.
\ifdefined\finalmode
  \renewcommand{\Added}[1]{#1}%
  \renewcommand{\Deleted}[1]{}%
  \renewcommand{\Replaced}[2]{#1}%
  \renewcommand{\added}[2][]{#2}%
  \renewcommand{\deleted}[2][]{}%
  \renewcommand{\replaced}[3][]{#2}%
  \renewenvironment{DeletedBlock}{\setbox0\vbox\bgroup}{\egroup}%
  \renewenvironment{AddedBlock}{\par}{\par}%
\fi

% Base16 Tomorrow Light colors
\definecolor{base00}{HTML}{ffffff} % Background
\definecolor{base01}{HTML}{e0e0e0} % Lighter background
\definecolor{base02}{HTML}{d6d6d6} % Selection
\definecolor{base03}{HTML}{8e908c} % Comments
\definecolor{base04}{HTML}{969896} % Dark foreground
\definecolor{base05}{HTML}{4d4d4c} % Default foreground
\definecolor{base06}{HTML}{282a2e} % Light foreground
\definecolor{base07}{HTML}{1d1f21} % Black
\definecolor{base08}{HTML}{c82829} % Red
\definecolor{base09}{HTML}{f5871f} % Orange
\definecolor{base0A}{HTML}{eab700} % Yellow
\definecolor{base0B}{HTML}{718c00} % Green
\definecolor{base0C}{HTML}{3e999f} % Cyan
\definecolor{base0D}{HTML}{4271ae} % Blue
\definecolor{base0E}{HTML}{8959a8} % Purple
\definecolor{base0F}{HTML}{a3685a} % Brown

\lstdefinestyle{base16python}{
  language=Python,
  morekeywords={pass} % doesn't seem to work..
  backgroundcolor=\color{base00},
  basicstyle={\renewcommand\ttdefault{cmtt}\ttfamily\small\color{base05}},
  keywordstyle=\color{base0E},
  commentstyle=\color{base03},
  stringstyle=\color{base0B},
  numberstyle=\scriptsize\color{base03},
  identifierstyle=\color{base05},
  showstringspaces=false,
  breaklines=true,
  postbreak=\space,
  frame=none,
  rulecolor=\color{base01},
  tabsize=4,
  columns=fullflexible,
  keepspaces=true,
  moredelim=[l][\color{base0B}]{@},
}

\lstnewenvironment{python}[1][]{\lstset{style=base16python,#1}}{}

% for use in inferrule
\newcommand{\pytext}[1]{%
  {\renewcommand\ttdefault{cmtt}\ttfamily\small\color{base05}#1}%
}
\newcommand{\pyinline}[1]{%
  \ifmmode
    \text{\pytext{#1}}%
  \else
    \mbox{\lstinline[style=base16python]{#1}}%
  \fi
}


\newtheorem{theorem}{Theorem}
\newtheorem{lemma}{Lemma}
\newtheorem{definition}{Definition}

\setlist[itemize]{leftmargin=*, nosep, topsep=2mm}
\setlist[enumerate]{nosep}

% PACMPL (PLDI) metadata; placeholders for a draft.
\acmJournal{PACMPL}
\acmVolume{10}
\acmNumber{PLDI}
\acmArticle{1}
\acmArticleSeq{0}
\acmYear{2026}
\acmMonth{6}
\setcopyright{none}
\settopmatter{printacmref=false}
\citestyle{acmauthoryear}

\begin{document}

\ifdefined\anonmode
  \title{A Pure Functional Subset of Python for Scientific Computing}
\else
  \title{\PurePy: A Functional Subset of Python for Scientific Computing}
\fi
\version{0}{14}{1}

\author{Tobias Antensteiner}
\affiliation{%
  \department{Institute of Theoretical Informatics}
  \institution{Karlsruhe Institute of Technology}
  \country{Germany}}

\author{Nikolaus Huber}
\affiliation{%
  \department{Department of Computer Science and Technology}
  \institution{University of Cambridge}
  \country{UK}}

\author{Dominic Orchard}
\affiliation{%
  \department{School of Computing}
  \institution{University of Kent}
  \country{UK}}
\affiliation{%
  \department{Department of Computer Science and Technology}
  \institution{University of Cambridge}
  \country{UK}}

\author{Jacob Pake}
\affiliation{%
  \department{School of Computing}
  \institution{University of Kent}
  \country{UK}}

\author{Roly Perera}
\affiliation{%
  \department{Department of Computer Science and Technology}
  \institution{University of Cambridge}
  \country{UK}}
\affiliation{%
  \department{School of Computer Science}
  \institution{University of Bristol}
  \country{UK}}

\renewcommand{\shortauthors}{Antensteiner et al.}

\begin{abstract}
\PurePy is a pure, side-effect-free subset of Python: a functional language whose programs are referentially
transparent, reproducible, and easy to reason about, yet remain ordinary Python. Every well-formed \PurePy
program is valid Python, and where \PurePy gives a run a result, Python gives the same one, so the same source
runs under a standard Python interpreter and under any conforming implementation. This paper specifies \PurePy as a versioned standard, with a
formal grammar, well-formedness rules based on a definite-assignment analysis, and an operational semantics
that agrees with Python on every run it defines. A reference checker decides well-formedness, and a conformance test suite is exercised on CPython, GraalPy, and Pyodide. Aimed first at researchers in programming languages and pedagogy, \PurePy
is intended to grow into a common language for scientific computing, supporting portable applications in
modelling, data processing, analysis, and visualisation.
\end{abstract}

\maketitle

\section{Introduction}

Python is among the most widely used languages in scientific computing, data analysis, and teaching, valued for
an accessible surface syntax and a large ecosystem. That reach comes with a large and permissive language.
Programs mutate state freely, depend on side effects and dynamic binding, and admit many implicit coercions, so
they can be hard to reason about, hard to reproduce, and hard to target with formal tools or with
implementations other than the reference interpreter. A researcher in programming languages or in pedagogy who
wants a well-behaved functional language faces a trade-off: such a language is easy to reason about but
unfamiliar, whereas Python is familiar but unruly.

\PurePy carves out a pure, side-effect-free subset of Python that is well-behaved and easy to reason about, yet
remains ordinary Python. Every well-formed \PurePy program is valid Python and runs under a standard
interpreter with the same result, so one artifact serves two roles: a program a scientist can run today, and a
common target that a research language or tool can produce and consume. Purity makes programs referentially
transparent, deterministic, and reproducible, the properties that scientific computing and formal reasoning
both want.

This paper defines \PurePy as a versioned formal standard: a formal grammar, well-formedness rules based on a
definite-assignment analysis, and an operational semantics that preserves the meaning of well-formed programs
in Python. A reference checker decides membership of the subset, and a conformance test suite exercises the
standard against CPython, PyPy, GraalPy, and Pyodide. The guarantee is precise: membership is decidable, and a
well-formed \PurePy program is valid Python whose \PurePy meaning is its Python meaning.

The design is deliberately conservative. Rather than reconstruct Python's full behaviour, \PurePy chooses a
simple, well-behaved core and excludes the idioms that would break it, among them mutable state, effectful
control, implicit coercions, and constructs whose meaning depends on module load order
(\secref{overview}). The standard is versioned, so the subset can grow as features earn their place
without disturbing programs written against an earlier version.

\PurePy is aimed first at researchers in programming languages and in pedagogy, and is intended to grow into a
common language for scientific computing, supporting portable applications in modelling, data processing,
analysis, and visualisation. Because every \PurePy program is also a Python program, adopting it costs nothing
at the surface, and a research language that targets \PurePy inherits a large and familiar runtime.

{\small
\begin{center}
\begin{minipage}[t]{0.4\textwidth}
\vspace{0pt}
\begin{tabular}{ll}
\textbf{To do} & \\
Reserved words & \issue{9} \\
Set and dict comprehensions & \issue{52} \\
Set and dict literals & \issue{87} \\
Membership operator (\pyinline{in}) & \issue{80} \\
Identity operator (\pyinline{is}) & \issue{81} \\
Chained comparison & \issue{82} \\
Enums & \issue{86} \\
Rules for remaining expression forms & \issue{133} \\
\end{tabular}
\end{minipage}\hspace{2em}
\begin{minipage}[t]{0.4\textwidth}
\vspace{0pt}
\begin{tabular}{ll}
\textbf{Under consideration} & \\
Destructuring assignment & \issue{54} \\
f-strings & \issue{55} \\
Default arguments & \issue{56} \\
\pyinline{*args} and \pyinline{**kwargs} & \issue{57} \\
Decorators & \issue{58} \\
Slicing & \issue{59} \\
Star patterns & \issue{84} \\
Or patterns & \issue{85} \\
Relative imports & \issue{126} \\
Function keyword arguments & \issue{127} \\
Postponed annotation evaluation & \issue{128} \\
Well-formedness correctness properties & \issue{130} \\
Primitive values and their application & \issue{132} \\
\end{tabular}
\end{minipage}
\end{center}}

\subsection{Contributions}

After setting out the problem \PurePy addresses in \secref{overview}, we make the following
contributions:
\begin{itemize}
\item a definition of \PurePy as a versioned standard, with a formal grammar and well-formedness rules based on
a definite-assignment analysis (\secref{syntax}) and an operational semantics that preserves the
meaning of well-formed programs in Python (\secref{opsem});
\item a reference checker that decides membership of the subset, validated by a conformance test suite
exercised on CPython, PyPy, GraalPy, and Pyodide (\secref{reference-impl});
\item conservative extensions to the core language, including partial application and piecewise function
definitions (\secref{extensions});
\item case studies that use \PurePy as a common functional language (\secref{case-studies});
\item an experimental evaluation of how well AI agents generate \PurePy and translate existing Python into it
(\secref{evaluation}).
\end{itemize}
\secref{related} summarises related efforts, and \secref{conclusion} presents our conclusions
and discusses future work.

\section{Overview}
\label{sec:overview}

\PurePy is a subset of Python chosen so that its programs are well-behaved and broadly \emph{pure}, yet run
unchanged under a standard Python interpreter and compute the same result. Purity here is informal: the
functional core has no observable side effects, and a name's meaning does not depend on evaluation order or on
mutation. A variable, once definitely assigned, denotes a fixed value; a function is fixed by its definition;
and a module's members are fixed once it has loaded. Python permits much that breaks this discipline, so
\PurePy tightens or excludes the idioms responsible. A precise characterisation of purity is left to future
work.

The exclusions fall into a few groups. Some features are imperative or effectful and have no place in a
functional core, such as loops, reassignment of captured variables, exceptions, and context managers. Others
are accepted by Python but give a program a meaning that depends on load order or runtime state, such as a name
that is at once an assignment and a submodule; \PurePy rejects these so that meaning is static. A third group is
excluded to keep the static account tractable, among them truthiness of non-boolean conditions and first-class
modules and classes. The rest of this section lists the excluded features and the Python runtime errors that
well-formedness rules out (\figref{prevented-errors}).

\subsection{Prohibited Python features}
\label{sec:prohibited}

The following Python features are excluded from \PurePy as they are incompatible with purity or
functional programming:
\begin{itemize}
\item \pyinline{while} loops, \pyinline{for} loops, \pyinline{break}, \pyinline{continue}
\item \pyinline{async}/\pyinline{await}
\item \pyinline{global}, \pyinline{nonlocal}
\item \pyinline{del}
\item \pyinline{try}/\pyinline{except}/\pyinline{raise}
\item Augmented assignment (\pyinline{+=}, \pyinline{-=}, etc.)
\item Walrus operator (\pyinline{:=})
\item \pyinline{yield}/\pyinline{yield from}
\item \pyinline{class} definitions other than dataclasses
\item \pyinline{with} statements
\item Truthiness: a non-boolean value used as a condition (the test of an \pyinline{if} or conditional expression, or an \pyinline{assert}). Python coerces any value to \pyinline{bool}; \PurePy requires the condition to be boolean.
\item List patterns matching tuple values, or tuple patterns matching list values. In \PurePy a list pattern matches only a list, and a tuple pattern matches only a tuple. Python permits either pattern against either kind of sequence.
\item Positional class patterns matching fewer sub-patterns than the class's field arity. \PurePy requires exact arity; Python permits any number from $0$ up to the class's arity.
\item Local or interleaved imports: an \pyinline{import} or \pyinline{from} statement inside a function or other block, or after a non-import statement. \PurePy permits imports only at the start of a module; Python allows them anywhere.
\item Wildcard imports (\pyinline{from q import *}); PEP 8 recommends avoiding these\footnote{\url{https://peps.python.org/pep-0008/}}.
\item A submodule importing itself through its parent (\pyinline{from a import b} inside module \pyinline{a.b}). \PurePy treats the from-imported submodule as a dependency, so this is an import cycle and is rejected; Python permits it, binding the partially-loaded module.
\item An import (not a from-import) naming a proper descendant of the importing module (\pyinline{import x.y} inside module \pyinline{x}): the binding for the head would denote the still-loading module. Python permits it, binding the partially-initialised module. A from-import of the module's own submodule is permitted.
\item First-class classes: in Python, class declarations are assignments, so classes flow as values. In \PurePy class names appear only at constructor and pattern positions, and declarations are module-level only.
\item Re-declaring a class name: a dataclass declaration whose name already denotes a class in the module. Python permits the re-declaration, leaving the earlier class reachable only through its existing instances; \PurePy requires a class name to be unique within its module, so that a qualified name identifies a class.
\item First-class modules: in Python, modules are values. In \PurePy a module name may be projected (member access) but not used as a value.
\item A module binding a name that is also the name of one of its submodules (\S\ref{sec:submodule-clash}).
\item Referring to a submodule the referring module does not itself import: \pyinline{import a} followed by \pyinline{a.b.f()} is rejected unless the module also imports \pyinline{a.b}. Python permits the reference whenever some earlier import, anywhere in the program, loaded \pyinline{a.b} (\S\ref{sec:uninitialised-members}).
\end{itemize}

Complementing the list above (programs Python runs that \PurePy rejects), a well-formed \PurePy
program avoids the Python runtime errors listed in \figref{prevented-errors}, under the conditions noted
there. Errors from ill-typed operations, such as wrong operand types, calling a non-function, or attribute
access on an object lacking the attribute, are outside the scope of well-formedness.

\subsubsection{Variable and submodule with the same name}
\label{sec:submodule-clash}

In Python, an attribute of a package can be bound by an assignment in the package's own body, and by a load
of the same-named submodule, which assigns the new module to the attribute. Suppose the body of \pyinline{pkg}
is the single assignment \pyinline{sub = 1}, and \pyinline{pkg} also has a submodule \pyinline{pkg.sub}:
\vspace{-\medskipamount}

\begin{minipage}[t]{0.48\textwidth}
\begin{python}
import pkg
print(pkg.sub)  # 1
\end{python}
\end{minipage}
\hfill
\begin{minipage}[t]{0.48\textwidth}
\begin{python}
import pkg.sub
print(pkg.sub)  # <module 'pkg.sub'>
\end{python}
\end{minipage}

On the left, \pyinline{pkg.sub} is never loaded, so the assignment is visible; on the right, the load
overwrites the assignment. The name \pyinline{pkg.sub} has no load-order-independent meaning, so the
module rules prohibit a module from binding a name that is also the name of one of its submodules.

\begin{figure}
\centering
\small
\begin{tabular}{>{\raggedright\arraybackslash}p{0.26\textwidth}>{\raggedright\arraybackslash}p{0.52\textwidth}>{\raggedright\arraybackslash}p{0.12\textwidth}}
\hline
\textbf{Python error} & \textbf{Ruled out when} & \textbf{See} \\
\hline
\pyinline{NameError}, \pyinline{UnboundLocalError} & Always: a variable may be referenced only where it is definitely assigned. & \figref{well-formed}, \figref{well-formed-expr} \\[1mm]
\pyinline{ModuleNotFoundError} & Always: an import requires its target to be a module of the program. & \figref{well-formed-imports} \\[1mm]
\pyinline{ImportError} & Always: a from-import requires each imported name to be a member of the source module, and import cycles are rejected. (A Python cycle may run or fail, depending on where the cycle is entered.) & \figref{imports}, \figref{well-formed-module} \\[1mm]
\pyinline{AttributeError} & Always: a module attribute reference must resolve to a definitely assigned member.  & \figref{well-formed-expr} \\[1mm]
\pyinline{TypeError} & The operation is a constructor call or class pattern, which must match the class's arity exactly. (Only excess sub-patterns are a Python error; a pattern with fewer is prohibited even though Python permits it.) & \figref{well-formed-expr}, \figref{well-formed-pat} \\
\hline
\end{tabular}
\caption{Python errors prevented by well-formedness}
\label{fig:prevented-errors}
\end{figure}

\section{Syntax}
\label{sec:syntax}

The syntax of \PurePy is given in \figref{syntax-stmt} and \figref{syntax-expr}. A \emph{block} is a non-empty sequence of
statements; either a top-level statement, a \pyinline{def} body, or a branch of a conditional. A
\emph{mutual region} $d_1 \ldots d_{n+1}$ is a maximal contiguous non-empty sequence of \pyinline{def}
statements in a block: any non-\pyinline{def} statement in that block, by definition, separates mutual
regions. Maximality is not enforced by the grammar in \figref{syntax-stmt}; the parser is responsible for
grouping consecutive \pyinline{def}s into a single region. The well-formedness rules treat each region as
one mutually-recursive binding (see \secref{mutual-recursion}).

\subsection{Lexical syntax}
\label{sec:lexical}

\todo{The grammar of \PurePy is an abstract syntax, and its metavariables range over values rather than
over source text: $n$ is a number and $w$ is a string, not the lexemes that denote them. How source text
becomes abstract syntax follows Python, including the indentation that delimits blocks; say where \PurePy
narrows those rules.}

\subsection{Contexts}
\label{sec:contexts}

Metavariables $x$, $f$, $c$ range over \emph{variables}; idiomatically $f$ is used when the variable is expected to denote a function and $c$ when it is expected to denote a class.

\begin{definition}[Context]
A context $\Gamma$ (or $\Delta$) is a finite map from variables to one of the following \emph{context entries} $\theta$:
\begin{itemize}
\item the variable's \emph{definite assignment} status $a \in \B = \set{\statusDefAssigned, \statusNotDefAssigned}$;
\item a \emph{module reference}, for a fully-qualified module name $q$: either a \emph{stub} $\modStub{q}$, or \emph{loaded} reference $\modLoaded{q}{\Gamma}$ carrying the \emph{signature} $\Gamma$ of $q$, the context specifying its members;
\item a \emph{class entry} $C$, of the form $\clsEntry{\Gamma}{q.c}{\seq{x}}{c'}$, recording the declaring context $\Gamma$, the qualified name $q.c$ of the class, sequence of distinct field names $\seq{x}$, and base class name $c'$ (or $\bot$ for no base).
\end{itemize}
\end{definition}

\noindent Module references arise from import statements. A module declares each class name at most once, so the qualified name of a class entry is the identity of the class. The definite-assignment terminology is borrowed from Java~\cite[\S16]{gosling2025}; $\statusDefAssigned$ indicates \emph{definitely assigned} and $\statusNotDefAssigned$ \emph{not definitely assigned} (or, equivalently, degenerate types ``any'' and ``none'').
It is useful to represent contexts as total functions on the set of all variables by extending the codomain with $\bot$ to indicate \emph{undefined}. We write $\Gamma, \bind{x}{\theta}$ for the context that sends $x$ to $\theta$ and otherwise agrees with $\Gamma$, define $\dom(\Gamma) = \set{x \mid \Gamma(x) \neq \bot}$, and write $\varnothing$ for the context where all variables are undefined. For a set of variables $X$, write $\statusDefAssigned(X)$ for the context $\set{\bind{x}{\statusDefAssigned} \mid x \in X}$; similarly $\statusNotDefAssigned(X)$.

\paragraph{Context operations.}

Contexts admit a non-commutative monoid of sequential composition $\Gamma \override \Delta$ (``override'') and
an idempotent commutative semigroup of parallel composition $\merge$ (``merge''). Together these form a
\emph{right near-semiring}~\cite{vanhoorn1967}: $\override$ right-distributes over $\merge$ (Lemma~1), but left
distributivity does not hold. Both operations are defined pointwise.
%
\begin{definition}[Context override]
{\small
\begin{align*}
\bot \override \theta &= \theta \override \bot = \theta \\
\theta \override \theta' &= \theta' \quad (\text{otherwise})
\end{align*}}
\end{definition}

\noindent Override is thus right-biased and has $\bot$ as both left and right unit.

\begin{definition}[Context merge]
For definite-assignment statuses $a, a' \in \B$:
{\small
\begin{align*}
\bot \merge \bot &= \bot \\
a \merge \bot = \bot \merge a &= \statusNotDefAssigned \\
a \merge a' &= a \meet a'
\end{align*}}
\end{definition}

\noindent Merge is otherwise undefined. Restricted to $\B$, it returns $\statusNotDefAssigned$ when exactly one argument is defined and the conjunction otherwise; so $\bot$ is not a unit.

\begin{definition}[Context extend]
\label{def:extend}
{\small
\begin{align*}
\modLoaded{q}{\Gamma} \extend \modLoaded{q}{\Gamma'} &= \modLoaded{q}{\Gamma \extend \Gamma'} \\
\modStub{q} \extend \modLoaded{q}{\Gamma} = \modLoaded{q}{\Gamma} \extend \modStub{q} &= \modLoaded{q}{\Gamma} \\
\bot \extend \theta = \theta \extend \bot &= \theta \\
\theta \extend \theta' &= \theta' \quad (\text{otherwise})
\end{align*}}
\end{definition}

\noindent Extend combines module references for the same module recursively, preferring loaded references to stubs, and is otherwise right-biased like override. It is used when processing import prefixes, where two imports rooted at the same package must combine rather than shadow.

\begin{lemma}[Right distributivity]
\[
(\Delta \merge \Delta') \override \Gamma = (\Delta \override \Gamma) \merge (\Delta' \override \Gamma).
\]
\end{lemma}

\noindent Left distributivity ($\Gamma \override (\Delta \merge \Delta') = (\Gamma \override \Delta) \merge
(\Gamma \override \Delta')$) does not hold. If it did, the following two programs would be equivalent, since
sequential composition with a preceding assignment could be distributed into the branches of a conditional.
But they are not: the left program is ill-formed, because a variable assigned before a conditional does not
remain definitely assigned if only some branches assign it (see also \secref{local-assignment}).
\vspace{-\medskipamount}

\begin{minipage}[t]{0.48\textwidth}
\begin{python}
def f(b):
    x = 1
    if b:
        x = 2
    else:
        pass
    return x  # error
\end{python}
\end{minipage}
\hfill
\begin{minipage}[t]{0.48\textwidth}
\begin{python}
def f(b):
    if b:
        x = 1
        x = 2
    else:
        x = 1
        pass
    return x  # ok
\end{python}
\end{minipage}

\begin{figure}
\small
\begin{center}
\begin{tabular}{rcll}
$s$ & $::=$ & $\exPass$ & (Pass) \\
    & $\mid$ & $\exAssign{x}{e}$ & (Assign) \\
    & $\mid$ & $e$ & (ExprStmt) \\
    & $\mid$ & $\exAssert{e}$ & (Assert) \\
    & $\mid$ & $\exAssertMsg{e}{e'}$ & (AssertMsg) \\
    & $\mid$ & $\exReturnNone$ & (ReturnNone) \\
    & $\mid$ & $\exReturn{e}$ & (Return) \\
    & $\mid$ & $\exIf{e}{b}{\exElif{e_1}{b_1}\;\ldots\;\exElif{e_n}{b_n}}$ & (If) \\
    & $\mid$ & $\exIfElse{e}{b}{\exElif{e_1}{b_1}\;\ldots\;\exElif{e_n}{b_n}}{b'}$ & (IfElse) \\
    & $\mid$ & $\exMatch{e}{\exCase{p_1}{b_1}\;\ldots\;\exCase{p_{n+1}}{b_{n+1}}}$ & (Match) \\
    & $\mid$ & $d_1 \;\ldots\; d_{n+1}$ & (MutualRegion) \\
    & $\mid$ & $\exDataclass{c}{\phi}$ & (Dataclass) \\
    & $\mid$ & $\exDataclassExtend{c}{c'}{\phi}$ & (DataclassExtend)
   \\[3mm]
$\iota$ & $::=$ & $\exImport{q}$ & (Import) \\
    & $\mid$ & $\exFromImport{q}{x_1, \ldots, x_{n+1}}$ & (FromImport)
   \\[3mm]
$\phi$ & $::=$ & $\exField{x_1} \;\ldots\; \exField{x_{n+1}}$ & (FieldList) \\
    & $\mid$ & $\pyinline{pass}$ & (EmptyFieldList)
   \\[3mm]
$b$ & $::=$ & $s_1 \;\ldots\; s_{n+1}$ & (Block)
   \\[3mm]
$t$ & $::=$ & $\iota_1 \;\ldots\; \iota_m \;\; s_1 \;\ldots\; s_n$ & (ModuleBody)
   \\[3mm]
$q$ & $::=$ & $x_1.x_2.\ldots.x_{n+1}$ & (Name)
   \\[3mm]
$d$ & $::=$ & $\exDef{f}{x_1, \ldots, x_n}{b}$ & (Def)
\end{tabular}
\end{center}
\caption{Modules and blocks}
\label{fig:syntax-stmt}
\end{figure}

\begin{figure}
\small
\centering
\ottgrammarfigure{
    \otte \ottinterrule
    \ottg \ottinterrule
    \ottl \ottinterrule
    \ottp \ottafterlastrule
}
\caption{Expressions and patterns}
\label{fig:syntax-expr}
\end{figure}

% \begin{center}
% \begin{tabular}{rcll}
% $e$ & $::=$ & $x$ & (Var) \\
%     & $\mid$ & $\ell$ & (Literal) \\
%     & $\mid$ & $\exApp{e}{e_1, \ldots, e_n}$ & (Call) \\
%     & $\mid$ & $\exConstruct{C}{e_1, \ldots, e_n}$ & (Construct) \\
%     & $\mid$ & $\exDict{x_1 \mathord{:} e_1, \ldots, x_n \mathord{:} e_n}$ & (Dict) \\
%     & $\mid$ & $\exProject{e}{x}$ & (Attribute) \\
%     & $\mid$ & $\exDProject{e}{e'}$ & (Subscript) \\
%     & $\mid$ & $\exBinaryApp{e}{\odot}{e'}$ & (BinOp) \\
%     & $\mid$ & $\exUnaryApp{\ominus}{e}$ & (UnaryOp) \\
%     & $\mid$ & $\exCondExpr{e_1}{e}{e_2}$ & (CondExpr) \\
%     & $\mid$ & $\exList{e_1, \ldots, e_n}$ & (List) \\
%     & $\mid$ & $\exTuple{e_1, \ldots, e_n}$ & (Tuple) \\
%     & $\mid$ & $\exListComp{e}{q_1\;\ldots\;q_n}$ & (ListComp) \\
%     & $\mid$ & $\exLambda{x_1, \ldots, x_n}{e}$ & (Lambda)
%    \\[3mm]
% $q$ & $::=$ & $\qualGuard{e}$ & (Guard) \\
%     & $\mid$ & $\qualGenerator{x}{e}$ & (Generator)
%    \\[3mm]
% $\ell$ & $::=$ & $n$ & (Number) \\
%     & $\mid$ & $\pyinline{True}$ & (True) \\
%     & $\mid$ & $\pyinline{False}$ & (False) \\
%     & $\mid$ & $\pyinline{None}$ & (None)
%    \\[3mm]
% $p$ & $::=$ & $\ell$ & (PatLiteral) \\
%     & $\mid$ & $-n$ & (PatNegLit) \\
%     & $\mid$ & $x$ & (PatVar) \\
%     & $\mid$ & $\patWild$ & (PatWildcard) \\
%     & $\mid$ & $\patList{p_1, \ldots, p_n}$ & (PatList) \\
%     & $\mid$ & $\patTuple{p_1, \ldots, p_n}$ & (PatTuple) \\
%     & $\mid$ & $\patClass{C}{p_1, \ldots, p_n, x_{n+1} \mathord{=} p_{n+1}, \ldots, x_{n+m} \mathord{=} p_{n+m}}$ & (PatClass) \\
%     & $\mid$ & $p\;\pyinline{as}\;x$ & (PatAs)
% \end{tabular}
% \end{center}


\paragraph{Operators.}

The metavariable $\odot$ ranges over the supported binary operators other than \pyinline{and} and
\pyinline{or}, and $\ominus$ over the supported unary
operators, both enumerated in \figref{precedence}. The operators \pyinline{and} and \pyinline{or} have their own
expression forms; they do not evaluate both operands. The grammar of \figref{syntax-expr} is intentionally ambiguous
on operator forms; disambiguation is by the precedence and associativity given in \figref{precedence},
which follow Python.

\begin{figure}
\small
\begin{center}
\begin{tabular}{rlll}
\hline
Level & Operators & Arity & Associativity \\
\hline
1 (lowest)  & \pyinline{or}                                                              & binary & left \\
2           & \pyinline{and}                                                             & binary & left \\
3           & \pyinline{not}                                                             & unary  & --- \\
4           & \pyinline{==}\quad\pyinline{!=}\quad\pyinline{<}\quad\pyinline{<=}\quad\pyinline{>}\quad\pyinline{>=} & binary & left \\
5           & \pyinline{+}\quad\pyinline{-}                                                & binary & left \\
6           & \pyinline{*}\quad\pyinline{/}\quad\pyinline{//}\quad\pyinline{\%}                & binary & left \\
7           & \pyinline{+}\quad\pyinline{-}                                                & unary  & --- \\
8 (highest) & \pyinline{**}                                                              & binary & right \\
\hline
\end{tabular}
\end{center}
\caption{Operator precedence (lowest binds loosest); follows Python.}
\label{fig:precedence}
\end{figure}


\subsection{Modules and programs}
\label{sec:modules-programs}

A \PurePy \emph{program} is a finite map $\mathcal{M}$ from \emph{fully-qualified module names}
(\figref{syntax-stmt}) to \emph{module bodies} $t$, of which one is a distinguished \emph{main module} named
$\pymath{\_\_main\_\_}$; here $\nameOf(q)$ is the string literal naming $q$, so
$\nameOf(\pymath{\_\_main\_\_}) = \pymath{"\_\_main\_\_"}$. We identify a module with its fully-qualified name, writing ``module $q$'' for the module named $q$. A module body is a sequence of imports
$\seq{\iota}$ followed by a sequence of statements $\seq{s}$; since imports are not statements, they can
appear neither in blocks (and hence not in function bodies) nor after the first non-import statement. The
well-formedness and evaluation rules check and evaluate the statements of each body $\mathcal{M}(q)$ with the
assignment $\exAssign{\pymath{\_\_name\_\_}}{\nameOf(q)}$ prepended, modelling the
$\pymath{\_\_name\_\_}$ binding the runtime supplies automatically. A second map $\mathcal{P}$ gives the
predefined modules of \figref{predefined}, which have no body: $\mathcal{P}$ takes each such name to its
environment (Definition~\ref{def:predefined}). The two domains are disjoint, and their union is
prefix-closed: the parent of any module in $\dom(\mathcal{M}) \cup \dom(\mathcal{P})$ is itself in
$\dom(\mathcal{M}) \cup \dom(\mathcal{P})$.\footnote{A source tree encodes such an $\mathcal{M}$: each module's body is its \texttt{.py}
file, each package's body is its \texttt{\_\_init\_\_.py}, and a package directory without an
\texttt{\_\_init\_\_.py} encodes the empty body (a namespace package, PEP 420, \url{https://peps.python.org/pep-0420/}). Prefix-closure thus holds by
construction.} The maps $\mathcal{M}$ and $\mathcal{P}$ are fixed and ambient in all
definitions that follow, as is the name $q$ of the enclosing module in the rules that check or evaluate a single module's body.


\begin{definition}[Predefined contexts and environments]
\label{def:predefined}
$\mathcal{P}$ takes each module $q$ of \figref{predefined} to an environment binding the members listed for
$q$, together with \pyinline{__name__}, to implementation-defined values. Its context is whatever types that
environment (\figref{well-formed-env}).
\end{definition}

\begin{figure}
\small
\centering
\begin{tabular}{lp{0.6\textwidth}}
\hline
\textbf{Module} & \textbf{Members} \\
\hline
\textup{\pyinline{builtins}} & \pyinline{print}, \pyinline{len}, \pyinline{range} \\
\pyinline{math} & \pyinline{pi}, \pyinline{e}, \pyinline{sqrt}, \pyinline{exp}, \pyinline{log}, \pyinline{sin}, \pyinline{cos}, \pyinline{tan}, \pyinline{floor}, \pyinline{ceil} \\
\pyinline{sys} & \pyinline{argv}, \pyinline{exit} \\
\pyinline{typing} & \pyinline{Any} \\
\pyinline{dataclasses} & \pyinline{dataclass} \\
\hline
\end{tabular}
\caption{Predefined modules and their members}
\label{fig:predefined}
\end{figure}

The metafunction $\assignsS$ extends by union from individual statements to blocks; on a module
body $\seq{\iota}\;\seq{s}$ it acts on $\seq{s}$. Write $q \prefixOf q'$ (\emph{prefix}) when $q' = q$ or $q' =
q.q''$ for some $q''$, and $\properPrefixOf$ for the strict order; the proper prefixes of a module name
are its \emph{ancestors}. The $\submods$ metafunction is
defined, with the class-related metafunctions, in \appendixref{sec:aux-modules}.

In this section, loading a module means checking its body and computing its signature
(\figref{well-formed-module}), the static counterpart of run-time loading (\secref{opsem}). Both
import forms load their target, together with its ancestors
(\figref{well-formed-imports}). A plain import also records each of these modules as a member of its
parent (the loads-as judgement), so that the binding created for the root of the target gives attribute
access along the whole chain. A from-import also loads any imported name that denotes a submodule. An
ancestor of the target is exempt from loading first if it is the importing module or one of its
ancestors, since at run time it may still be loading. The exemption applies only to from-imports, because
for a plain import
the only ancestor that could never load first is the importing module itself, and an import naming a proper
descendant of the importing module is excluded outright (\secref{no-load-dependent}). Loading a module
may thus require other modules to load first; the rules are interpreted inductively, so a program with an
import cycle is ill-formed, because a module in a cycle could finish loading only after itself. Program
well-formedness (\figref{well-formed-module}) requires the main module to be well-formed and hence,
recursively, every module the program can load; a module that is never imported is not checked.

\subsection{Well-formedness}

\subsubsection{Auxiliary definitions}

Three syntactic analyses support the well-formedness rules: $\assignsS(s)$, the set of variables assigned
anywhere in $s$ (without descending into nested function definitions, which have their own scope, so that
$\assignsS(s) = \dom(\Delta)$ when $s : \Assigns{\Delta}$); $\binds(p)$, the set of variables introduced by a
pattern $p$; and $\capturesS(s)$, the set of variables from the enclosing scope that are captured by closures
(lambdas or nested function definitions) within $s$, with companion $\capturesE(e)$ on expressions. All are
defined, by their principal cases, in \appendixref{sec:aux-syntactic}. $\capturesS$ extends to blocks and
module bodies by union (like $\assignsS$), with the singleton block identified with its statement.

\begin{definition}[Pattern subsumption]
The preorder $p \subsumed p'$ (\appendixref{sec:aux-subsumption}) holds when every value matched by $p$ is also matched by $p'$.
\end{definition}

\noindent Subsumption is used by the well-formedness rule for $\pyinline{match}$ to rule out unreachable cases and to check exhaustiveness; with the current pattern set, the latter reduces to requiring a catch-all (var or wildcard) case.

\subsubsection{Well-formedness rules}

A pattern list is well-formed (\figref{well-formed-pat}) when each pattern is linear (no repeated variable names) and no earlier pattern subsumes a later one. The \pyinline{match} well-formedness rules use this judgement to enforce reachability and linearity of case patterns.

\begin{figure}
\flushleft \shadebox{$\Gamma \vdash p_1 \cons \ldots \cons p_n$}
\begin{mathpar}
\small
\inferrule*[lab={\ruleName{pat-lit}}]
{ \strut }
{ \Gamma \vdash \ell }
\and
\inferrule*[lab={\ruleName{pat-neg-lit}}]
{ \strut }
{ \Gamma \vdash -n }
\and
\inferrule*[lab={\ruleName{pat-var}}]
{ \strut }
{ \Gamma \vdash x }
\and
\inferrule*[lab={\ruleName{pat-wild}}]
{ \strut }
{ \Gamma \vdash \patWild }
\and
\inferrule*[lab={\ruleName{pat-list}}]
{ \Gamma \vdash p_i \quad (\forall i) }
{ \Gamma \vdash \patList{p_1, \ldots, p_n} }
\and
\inferrule*[lab={\ruleName{pat-tuple}}]
{ \Gamma \vdash p_i \quad (\forall i) }
{ \Gamma \vdash \patTuple{p_1, \ldots, p_n} }
\and
\inferrule*[lab={\ruleName{pat-as}}]
{ \Gamma \vdash p }
{ \Gamma \vdash p\;\pyinline{as}\;x }
\and
\inferrule*[lab={\ruleName{pat-class}}]
{
  \Gamma \vdash q \names d(x_1, \ldots, x_n, x_{n+1}, \ldots, x_{n+m}) \\
  \set{y_1, \ldots, y_m} = \set{x_{n+1}, \ldots, x_{n+m}} \\
  \Gamma \vdash p_i \quad (\forall i)
}
{
  \Gamma \vdash \patClass{q}{p_1, \ldots, p_n, y_1 \mathord{=} p_{n+1}, \ldots, y_m \mathord{=} p_{n+m}}
}
\and
\inferrule*[lab={\ruleName{cases-nil}}]
{ \strut }
{ \Gamma \vdash \varepsilon }
\and
\inferrule*[lab={\ruleName{cases-cons}}]
{
  \Gamma \vdash p_1 \\
  \Gamma \vdash p_2 \cons \ldots \cons p_n \\
  \text{variables in } p_1 \text{ distinct} \\
  p_i \not\subsumed p_1 \quad (\forall i.\, 2 \le i \le n)
}
{
  \Gamma \vdash p_1 \cons p_2 \cons \ldots \cons p_n
}
\end{mathpar}
\caption{Case patterns are well-formed under $\Gamma$}
\label{fig:well-formed-pat}
\end{figure}


Every well-formed statement either \emph{definitely returns}, or \emph{definitely assigns} a context
$\Delta$, which may be empty. The base cases are that \pyinline{return e} definitely returns, and \pyinline{x
= e} definitely assigns $\set{\bind{\pyinline{x}}{\statusDefAssigned}}$. \emph{Result types} $R$ classify these outcomes: $R$ is either $\Returns$ (definitely returns) or $\Assigns{\Delta}$ (definitely assigns $\Delta$). Block well-formedness is given in
\figref{well-formed}; expression well-formedness in \figref{well-formed-expr}.

Context operations $\override$ and $\merge$ lift to a right semiring over result types $R$, with $\Assigns{\varnothing}$ as left and right unit for $\override$ and $\Returns$ as unit for $\merge$ and zero for $\override$.

\begin{definition}[Override and merge on result types]
{\small
\begin{align*}
\Assigns{\Delta} \merge \Assigns{\Delta'} & = \Assigns{(\Delta \merge \Delta')} \\
R \merge \Returns = \Returns \merge R & = R \\
\\
\Assigns{\Delta} \override \Assigns{\Delta'} & = \Assigns{(\Delta \override \Delta')} \\
R \override \Returns = \Returns \override R & = \Returns
\end{align*}}
\end{definition}

\noindent The $\Returns \override R$ case never arises because the well-formedness judgement excludes unreachable statements (\secref{unreachable}).

\begin{lemma}[Right distributivity (result types)]
\[
(R \merge R') \override S = (R \override S) \merge (R' \override S)
\]
\end{lemma}

\noindent A module attribute reference is well-formed only when it resolves to a definitely assigned member
(rule \ruleName{attr-module}); see \secref{no-load-dependent}.

\begin{figure}
\flushleft \shadebox{$\Gamma; \Lambda \vdash b : R$}
\begin{mathpar}
\small
\inferrule*[lab={\ruleName{pass}}]
{
  \strut
}
{
  \Gamma; \Lambda \vdash \exPass : \Assigns{\varnothing}
}
\and
\inferrule*[lab={\ruleName{assign}}]
{
  \Gamma; \Lambda \vdash e \\
  x \notin \capturesE(e)
}
{
  \Gamma; \Lambda \vdash \exAssign{x}{e} : \Assigns{\set{x : \statusDefAssigned}}
}
\and
\inferrule*[lab={\ruleName{expr-stmt}}]
{
  \Gamma; \Lambda \vdash e
}
{
  \Gamma; \Lambda \vdash e : \Assigns{\varnothing}
}
\and
\inferrule*[lab={\ruleName{assert}}]
{
  \Gamma; \Lambda \vdash e
}
{
  \Gamma; \Lambda \vdash \exAssert{e} : \Assigns{\varnothing}
}
\and
\inferrule*[lab={\ruleName{assert-msg}}]
{
  \Gamma; \Lambda \vdash e \\
  \Gamma; \Lambda \vdash e'
}
{
  \Gamma; \Lambda \vdash \exAssertMsg{e}{e'} : \Assigns{\varnothing}
}
\and
\inferrule*[lab={\ruleName{return-none}}]
{
  \strut
}
{
  \Gamma; \Lambda \vdash \exReturnNone : \Returns
}
\and
\inferrule*[lab={\ruleName{return}}]
{
  \Gamma; \Lambda \vdash e
}
{
  \Gamma; \Lambda \vdash \exReturn{e} : \Returns
}
\and
\inferrule*[lab={\ruleName{if}}]
{
  \Gamma; \Lambda \vdash e_i \quad (\forall i) \\
  \Gamma; \Lambda \vdash b_i : R_i \quad (\forall i)
}
{
  \Gamma; \Lambda \vdash \exIf{e_1}{b_1}{\exElif{e_2}{b_2}\;\ldots\;\exElif{e_{n+1}}{b_{n+1}}} : R_1 \merge \ldots R_{n + 1} \merge \ldots \merge \Assigns{\varnothing}
}
\and
\inferrule*[lab={\ruleName{if-else}}]
{
  \Gamma; \Lambda \vdash e_i \quad (\forall i) \\
  \Gamma; \Lambda \vdash b_i : R_i \quad (\forall i)
}
{
  \Gamma; \Lambda \vdash \exIfElse{e_1}{b_1}{\exElif{e_2}{b_2}\;\ldots\;\exElif{e_{n+1}}{b_{n+1}}}{b_{n+2}} : R_1 \merge \ldots \merge R_{n+2}
}
\and
\inferrule*[lab={\ruleName{mutual}}]
{
  \Gamma \runion \statusDefAssigned(\seq{f}) \runion \statusDefAssigned(\seq{x}_{i}) \runion
  \statusNotDefAssigned(\assignsS(b_i) \setminus \seq{x}_i); \Lambda \vdash b_i : R_i
  \quad (\forall i)
  \\
  f_1, \ldots, f_{n+1} \text{ distinct}
}
{
  \Gamma; \Lambda \vdash \exDef{f_1}{\vec{x}_1}{b_1} \cons \ldots \cons \exDef{f_{n+1}}{\vec{x}_{n+1}}{b_{n+1}} : \Assigns{\statusDefAssigned(\seq{f})}
}
\and
\inferrule*[lab={\ruleName{match}}]
{
  \Gamma; \Lambda \vdash e \\
  \Gamma \runion \statusDefAssigned(\binds(p_i)); \Lambda \vdash b_i : R_i \quad (\forall i) \\
  S_i = \Assigns{\statusDefAssigned(\binds(p_i))} \runion R_i \quad (\forall i) \\
  \Lambda \vdash p_1 \cons \ldots \cons p_{n+1} \\
  p_{n+1} \in \set{x, \patWild}
}
{
  \Gamma; \Lambda \vdash \exMatch{e}{\exCase{p_1}{b_1}\;\ldots\;\exCase{p_{n+1}}{b_{n+1}}} : S_1 \merge \ldots \merge S_{n+1}
}
\and
\inferrule*[lab={\ruleName{match-partial}}]
{
  \Gamma; \Lambda \vdash e \\
  \Gamma \runion \statusDefAssigned(\binds(p_i)); \Lambda \vdash b_i : R_i \quad (\forall i) \\
  S_i = \Assigns{\statusDefAssigned(\binds(p_i))} \runion R_i \quad (\forall i) \\
  \Lambda \vdash p_1 \cons \ldots \cons p_{n+1} \\
  p_{n+1} \notin \set{x, \patWild}
}
{
  \Gamma; \Lambda \vdash \exMatch{e}{\exCase{p_1}{b_1}\;\ldots\;\exCase{p_{n+1}}{b_{n+1}}} : S_1 \merge \ldots \merge S_{n+1} \merge \Assigns{\varnothing}
}
\and
\inferrule*[lab={\ruleName{cons}}]
{
  \Gamma; \Lambda \vdash s : \Assigns{\Delta} \\
  \Gamma \runion \Delta; \Lambda \vdash b : R \\
  \capturesS(s) \cap \assignsS(b) = \varnothing
}
{
  \Gamma; \Lambda \vdash s \cons b : \Assigns{\Delta} \runion R
}
\end{mathpar}
\caption{Block $b$ is well-formed in $\Gamma$ and $\Lambda$}
\label{fig:well-formed}
\end{figure}

\begin{figure}
\begin{subfigure}{\textwidth}
\flushleft \shadebox{$\Gamma \vdash q \names \theta$}
\begin{mathpar}
\small
\inferrule*[lab={\ruleName{simple-module}}]
{
  x : q \in \Gamma
}
{
  \Gamma \vdash x \names q
}
\and
\inferrule*[lab={\ruleName{simple-class}}]
{
  c : (q', \vec{y}, c') \in \Gamma \\
  \fields(q'.c) = \vec{x}
}
{
  \Gamma \vdash c \names q'.c(\vec{x})
}
\and
\inferrule*[lab={\ruleName{qualified-module}}]
{
  \Gamma \vdash q \names q' \\
  q'.x \in \dom(\mathcal{M})
}
{
  \Gamma \vdash q.x \names q'.x
}
\and
\inferrule*[lab={\ruleName{qualified-class}}]
{
  \Gamma \vdash q \names q' \\
  \Gamma_{q'}(c) = (q', \vec{y}, c'') \\
  \fields(q'.c) = \vec{x}
}
{
  \Gamma \vdash q.c \names q'.c(\vec{x})
}
\end{mathpar}
\subcaption{Name $q$ resolves to context entry $\theta$ under $\Gamma$}
\end{subfigure}

\vspace{2mm}

\begin{subfigure}{\textwidth}
\flushleft \shadebox{$\Gamma \vdash e$}
\begin{mathpar}
\small
\inferrule*[lab={\ruleName{var}}]
{
  \Gamma(x) = \statusDefAssigned
}
{
  \Gamma \vdash x
}
\and
\inferrule*[lab={\ruleName{lambda}}]
{
  \Gamma \runion \statusDefAssigned(\seq{x}) \vdash e
}
{
  \Gamma \vdash \exLambda{\seq{x}}{e}
}
\and
\inferrule*[lab={\ruleName{constr}}]
{
  \Gamma \vdash q \names d(x_1, \ldots, x_n) \\
  \Gamma \vdash e_i \quad (\forall i)
}
{
  \Gamma \vdash \exConstruct{q}{e_1, \ldots, e_n}
}
\and
\inferrule*[lab={\ruleName{attr-module}}]
{
  \Gamma \vdash e \names q \\
  y \in \dom(\Gamma_q)
}
{
  \Gamma \vdash e.y
}
\and
\inferrule*[lab={\ruleName{attr-object}}]
{
  \Gamma \vdash e
}
{
  \Gamma \vdash e.y
}
\end{mathpar}
\subcaption{Expression $e$ is well-formed under $\Gamma$}
\end{subfigure}
\caption{Expression well-formedness}
\label{fig:well-formed-expr}
\end{figure}

\begin{figure}
\flushleft \shadebox{$\Gamma \vdash m : \Delta$}
\begin{mathpar}
\small
\inferrule*[lab={\ruleName{empty}}]
{
  \strut
}
{
  \Gamma \vdash \varepsilon : \varnothing
}
\and
\inferrule*[lab={\ruleName{import}}]
{
  M' = x_1.x_2.\ldots.x_{n+1} \\
  \Gamma \runion \statusDefAssigned(\set{x_1}) \vdash m : \Delta
}
{
  \Gamma \vdash \exImport{M'} \cons m : \statusDefAssigned(\set{x_1}) \runion \Delta
}
\and
\inferrule*[lab={\ruleName{from-import}}]
{
  \Gamma \runion \statusDefAssigned(\set{x_1, \ldots, x_{n+1}}) \vdash m : \Delta
}
{
  \Gamma \vdash \exFromImport{M'}{x_1, \ldots, x_{n+1}} \cons m : \statusDefAssigned(\set{x_1, \ldots, x_{n+1}}) \runion \Delta
}
\and
\inferrule*[lab={\ruleName{statement}}]
{
  \Gamma \vdash s : \Assigns{\Delta} \\
  \Gamma \runion \Delta \vdash m : \Delta' \\
  \capturesS(s) \cap \assignsS(m) = \varnothing
}
{
  \Gamma \vdash s \cons m : \Delta \runion \Delta'
}
\and
\inferrule*[lab={\ruleName{class}}]
{
  \decls(\phi) = x_1, \ldots, x_n \\
  x_1, \ldots, x_n \text{ distinct} \\
  \Gamma \vdash m : \Delta
}
{
  \Gamma \vdash \exDataclass{C}{\phi} \cons m : \Delta
}
\and
\inferrule*[lab={\ruleName{class-extend}}]
{
  B \in \dom(\Lambda_M) \\
  \decls(\phi) = x_1, \ldots, x_n \\
  x_1, \ldots, x_n \text{ distinct and disjoint from } \fields(M.B) \\
  \Gamma \vdash m : \Delta
}
{
  \Gamma \vdash \exDataclassExtend{C}{B}{\phi} \cons m : \Delta
}
\end{mathpar}
\caption{Body $m$ of module $M$ is well-formed in $\Gamma$}
\label{fig:well-formed-module}
\end{figure}


The signature $\Gamma$ a module's load judgement produces (\figref{well-formed-module}) comprises the classes
and variables a module defines, together with its intrinsic \pyinline{__name__} and a stub for each submodule,
never the names it imports: following PEP 484\footnote{\url{https://peps.python.org/pep-0484/}}, an imported
name is a private dependency, not re-exported as part of the interface. Submodule stubs and the module's own
bindings are disjoint, because the module rules exclude a module that binds a name that is also the name of
a submodule (\secref{no-load-dependent}).

\subsection{Examples}

\subsubsection{Assignment as definition}
\label{sec:assignment-as-definition}

Although \PurePy is pure and all variable usage is well-scoped, local variables are introduced using an
\emph{assignment statement}, for compatibility with Python. If a variable is definitely assigned in any branch
of a conditional, then for that variable to be usable after the conditional, every other branch of the
conditional must either definitely return, or definitely assign that variable too. Thus the following are both
well-formed:

\begin{figure}[!ht]
\centering
\begin{subfigure}{0.48\textwidth}
\centering
\begin{python}
if a < b:
   comp = "smaller"
else:
   comp = "bigger"
return comp
\end{python}
\end{subfigure}
\hfill
\begin{subfigure}{0.48\textwidth}
\centering
\begin{python}
if a < b:
   return "smaller"
else:
   comp = "bigger"
return comp
\end{python}
\end{subfigure}
\end{figure}

\noindent The definite assignment set for a conditional that has at least one branch which definitely assigns
is calculated by merging the definite assignment sets for all branches which definitely assign, using the
$\merge$ operator on environments.

Definite assignment thus gives each reference a unique lexical referent: a variable is introduced by
an assignment, possibly jointly by the branches of a conditional, and a name may be referenced only where one
such introduction covers every path. This is stricter than avoiding Python's \pyinline{NameError}: the
left-hand program of \secref{contexts} assigns \pyinline{x} on every run, but is rejected because the
reference would denote the branch-local variable on one path and the variable it hides on the other.

\subsubsection{Local assignment in conditional clauses}
\label{sec:local-assignment}

For flexibility, we permit \emph{branch-local} variables. The assignment of \pyinline{x} below (left) is
permitted, and can be used locally in that branch of the conditional:

\begin{figure}[!ht]
\centering
\begin{subfigure}{0.48\textwidth}
\centering
\begin{python}
def foo():
  z = 6
  if b:
    x = 3
    y = x + 1
  else:
    y = 0
  return z
\end{python}
\caption{Branch-local definition of \pyinline{x}}
\end{subfigure}
\hfill
\begin{subfigure}{0.48\textwidth}
\centering
\begin{python}
def foo():
  x = 6
  if b:
    x = 3
    y = x + 1
  else:
    y = 0
  return x # error: partially defined
\end{python}
\caption{Now with use of \pyinline{x} after the conditional}
\end{subfigure}
\end{figure}

\noindent However such variables only appear to be branch-local. The definition of $\merge$ actually ensures
that a variable assigned in any branch is considered assigned by the conditional; however unless it is
assigned in \emph{every} branch that does not definitely return, it will be bound to $\statusNotDefAssigned$ in the
corresponding environment. Any attempt to refer to such a variable is illegal, as it refers to a variable
which may not be assigned. (Thus the right-hand program is ill-formed, despite the enclosing \pyinline{x = 6},
because in the \pyinline{return x} statement, the inner assignment shadows the the outer one, and the inner
assignment is only partial.) If we were to allow the program on the right, there would be no straightforward
way to assign it a semantics which is both pure, and in agreement with Python, since in Python the function
returns \pyinline{3} in the case where \pyinline{b} is \pyinline{True}.

\subsubsection{Function-level scoping of local variables}

In Python, all local variables are scoped to the entire function body, not just from the point of
assignment onwards. The $\assignsS$ function in the \ruleName{def} rule reflects this: all variables
assigned anywhere in the function body are introduced at $\statusNotDefAssigned$ before the body is checked. This means
the following program is statically rejected by \PurePy, whereas Python rejects it at runtime with an
\pyinline{UnboundLocalError}:

\begin{python}
y = 7

def f():
    x = y  # error: y is local (due to assignment below)
    y = 8

f()
\end{python}

\noindent Because \pyinline{y} appears on the left-hand side of an assignment in \pyinline{f}, it is in
$\assignsS(b)$ for the body of \pyinline{f}, and is therefore introduced at $\statusNotDefAssigned$ at the top of the
function. The reference to \pyinline{y} in \pyinline{x = y} then fails the well-formedness check, since
$\Gamma(\pyinline{y}) = \statusNotDefAssigned$ rather than $\Gamma(\pyinline{y}) = \statusDefAssigned$.

\subsubsection{Reassignment of captured variables}
\label{sec:captured-vars}

The exclusion of captured-variable reassignment (\secref{no-mutable-state}) is enforced by a side condition
on the \ruleName{cons} rule: no variable captured by a statement $s$ may be reassigned in the rest of the block
$b$. In the example of \secref{no-mutable-state}, \pyinline{x} is captured by the definition of
\pyinline{g} and reassigned by \pyinline{x = 6}, so the program is rejected. A dataclass declaration counts as
an assignment of its class name for these definitions, though re-declaring a class name is excluded outright
(\secref{no-first-class}).

The same issue arises when a variable is captured and reassigned in the \emph{same} statement:

\begin{python}
x = 5
def x():
    return x
x()  # Python: returns the function (late binding); pure: returns 5
\end{python}

\noindent The side condition on the \ruleName{assign} rule rejects programs like this, requiring that the
variable being bound is not captured by the expression. The \ruleName{def} rule, by contrast, binds
each $f_i$ in the typing context for every body $b_i$, so the function names may legitimately appear in
$\capturesS$ of those bodies. This is what permits self-recursion (the singleton-region case) and
mutual recursion within a region.

\subsubsection{Mutual recursion}
\label{sec:mutual-recursion}

The \ruleName{def} rule binds the names $f_1, \ldots, f_{n+1}$ of a mutual region $d_1 \ldots d_{n+1}$
simultaneously. Each body $b_i$ is checked in a context that includes all the $f_j$ as definitely
assigned, so any $f_j$ may reference any other $f_k$ within the region (the singleton case $n = 0$
permits self-recursion). The names in a region must be distinct.

Inserting any non-\pyinline{def} statement between two \pyinline{def}s splits them into separate regions,
breaking the mutual reference:

\begin{python}
def even(n):
    if n == 0:
        return True
    return odd(n - 1)

x = even(0)  # intervening statement separates the regions

def odd(n):
    if n == 0:
        return False
    return even(n - 1)
\end{python}

\noindent This program is rejected statically: \pyinline{even} forms a singleton region in which
\pyinline{odd} is not yet bound, so the reference to \pyinline{odd} in its body fails the
\ruleName{var} rule. Python accepts it: the call \pyinline{even(0)} does not actually reach
\pyinline{odd}, and Python's late binding means \pyinline{odd} is only looked up when needed.

Each branch of a conditional may contain its own mutual regions, subject to the usual definite assignment
rules. The following is well-formed because both branches define both names:

\begin{python}
b = True

if b:
    def f():
        return g()
    def g():
        return "via mutual region"
else:
    def f():
        return "f independent"
    def g():
        return "g independent"

print(f())
print(g())
\end{python}

A region containing two definitions of the same name is rejected directly by the \ruleName{def}
rule's distinct-names side condition:

\begin{python}
def g():
    return 0

def f():
    return g()
def g():     # error: g defined twice in the same region
    return 1
\end{python}

\noindent (Python accepts this; late binding makes \pyinline{f}'s call to \pyinline{g} resolve to the
second \pyinline{def g}.)

Wrapping the second \pyinline{def f; def g} pair in a conditional sidesteps the duplicate-name issue:
the outer \pyinline{def g} and the inner \pyinline{def g} are now in different blocks, so they no longer
fall in the same region. One might still expect a captured-variable problem, since the inner
\pyinline{def g} could be seen as rebinding the outer \pyinline{g} that \pyinline{f} captures. No such
problem arises:

\begin{python}
def g():
    return 0

if b:
    def f():
        return g()
    def g():
        return 1
\end{python}

\noindent The inner \pyinline{def f; def g} form a branch-local mutual region, so \pyinline{f}'s body
resolves \pyinline{g} to its sibling rather than to the outer \pyinline{g}. The two are bound
simultaneously, so \pyinline{f} cannot be called before the inner \pyinline{g} exists, and the value
\pyinline{f} sees is the same as the value Python's late binding finds at call time.

\subsubsection{Pattern variables escape the match}

A variable bound by a \pyinline{case} pattern is in scope after the \pyinline{match} statement, just as
a variable assigned inside a conditional branch is in scope after the conditional. Python treats matching
this way; \PurePy follows.

\begin{python}
v = 5
match v:
    case x:
        pass
print(x)    # prints 5
\end{python}

Here the catch-all \pyinline{case x:} binds \pyinline{x} in the body's scope, and the binding survives into
the enclosing scope. Definite assignment after the match follows the same merge logic as a conditional: if
every case binds the variable (or returns), it is definitely assigned afterwards; otherwise it is not (see
also \secref{local-assignment}).

\subsubsection{Unreachable code}
\label{sec:unreachable}
A non-empty sequence of statements returns only if the last statement returns, and all preceding statements
assign. Thus the following is illegal:

\begin{python}
def foo(x):
  return x + 1
  return x    # error: unreachable
\end{python}

\begin{DeletedBlock}
\subsubsection{Conditionals without \pyinline{else} clauses}

An \pyinline{if} statement without an \pyinline{else} clause is equivalent to an \pyinline{if} statement with
an \pyinline{else} clause whose body is a $\exPass$ statement; such a statement never returns.

\subsubsection{Implicit \pyinline{None}}

A \pyinline{return} statement without an expression is equivalent to \pyinline{return None}.

\subsubsection{Implicit \pyinline{return}}

A function whose body $s$ assigns rather than returns is equivalent to a function whose body is $s$ followed
by \pyinline{return None}. This ensures every function returns, even if it is not explicitly specified.

\subsubsection{\pyinline{assert} without a message}

An \pyinline{assert e} statement without a message is equivalent to \pyinline{assert e, ""}.
\end{DeletedBlock}

\section{Well-formedness}
\label{sec:well-formedness}

\Added{Well-formedness is a definite-assignment analysis over the syntax of \secref{syntax}: it decides which
programs \PurePy gives a meaning to. The judgements are stated over contexts, and the section follows the
syntax, taking patterns, then blocks and expressions, then modules and programs. \secref{opsem} gives each
of these a dynamic counterpart.}

\subsection{Auxiliary definitions}

Three syntactic analyses support the well-formedness rules: $\assignsS(s)$, the set of variables assigned
anywhere in $s$ (without descending into nested function definitions, which have their own scope, so that
$\assignsS(s) = \dom(\Delta)$ when $s : \Assigns{\Delta}$); $\binds(p)$, the set of variables introduced by a
pattern $p$; and $\capturesS(s)$, the set of variables from the enclosing scope that are captured by closures
(lambdas or nested function definitions) within $s$, with companion $\capturesE(e)$ on expressions. All are
defined, by their principal cases, in \Replaced{\figref{aux}}{\appendixref{sec:aux-syntactic}}. $\capturesS$ extends to blocks and
module bodies by union (like $\assignsS$), with the singleton block identified with its statement.
\Added{There $\fvE(e)$ and $\fvS(b)$ are the free variables of an expression and of a block, defined in the
usual way. On the cases not shown, $\binds$ and $\capturesE$ recurse structurally, taking the union over
immediate subterms, and $\capturesS$ takes the union of $\capturesE$ over its subexpressions and $\capturesS$
over its subblocks. $\assignsS$ is empty on the statements not shown, is the union over the branch blocks of
a conditional, and does not descend into nested function definitions, which have their own scope.}

\begin{figure}
\small
\begin{align*}
\assignsS(\exAssign{x}{e}) &= \set{x} \\
\assignsS(\exDef{f_1}{\seq{x}_1}{b_1} \cons \ldots \cons \exDef{f_{n+1}}{\seq{x}_{n+1}}{b_{n+1}}) &= \set{f_1, \ldots, f_{n+1}} \\
\assignsS(\exMatch{e}{\exCase{p_1}{b_1}\;\ldots\;\exCase{p_{n+1}}{b_{n+1}}}) &= \textstyle\bigcup_i \big(\binds(p_i) \cup \assignsS(b_i)\big) \\
\assignsS(\exDataclass{c}{\phi}) &= \set{c} \\
\assignsS(\exDataclassExtend{c}{c'}{\phi}) &= \set{c} \\
\\
\binds(x) &= \set{x} \\
\binds(p\;\pyinline{as}\;x) &= \binds(p) \cup \set{x} \\
\\
\capturesS(\exDef{f_1}{\seq{x}_1}{b_1} \cons \ldots \cons \exDef{f_{n+1}}{\seq{x}_{n+1}}{b_{n+1}}) &= \textstyle\bigcup_i \big(\fvS(b_i) \setminus (\set{\seq{x}_i} \cup \assignsS(b_i))\big) \setminus \set{f_1, \ldots, f_{n+1}} \\
\capturesS(\exMatch{e}{\exCase{p_1}{b_1}\;\ldots\;\exCase{p_{n+1}}{b_{n+1}}}) &= \capturesE(e) \cup \textstyle\bigcup_i \big(\capturesS(b_i) \setminus \binds(p_i)\big) \\
\\
\capturesE(\exLambda{\seq{x}}{e}) &= \fvE(e) \setminus \set{\seq{x}}
\end{align*}
\caption{Assigned, bound, and captured variables: principal cases}
\label{fig:aux}
\end{figure}


\begin{definition}[Pattern subsumption]
The preorder $\Replaced{\Gamma \vdash p \subsumed p'}{p \subsumed p'}$ (\Replaced{\figref{subsumption}}{\appendixref{sec:aux-subsumption}}) holds when every value matched by $p$ is also matched by $p'$.
\end{definition}

\noindent Subsumption is used by the well-formedness rule for $\pyinline{match}$ to rule out unreachable cases and to check exhaustiveness; with the current pattern set, the latter reduces to requiring a catch-all (var or wildcard) case.

\begin{figure}
\flushleft \shadebox{$\added{\Gamma} \mathrel{\added{\vdash}} p \subsumed p'$}
\begin{mathpar}
\small
\inferrule*[lab={\ruleName{sub-lit}}]
{
  \strut
}
{
  \added{\Gamma} \mathrel{\added{\vdash}} \ell \subsumed \ell
}
\and
\inferrule*[lab={\ruleName{sub-neg-lit}}]
{
  \strut
}
{
  \added{\Gamma} \mathrel{\added{\vdash}} -n \subsumed -n
}
\and
\inferrule*[lab={\ruleName{sub-var}}]
{
  \strut
}
{
  \added{\Gamma} \mathrel{\added{\vdash}} p \subsumed x
}
\and
\inferrule*[lab={\ruleName{sub-wild}}]
{
  \strut
}
{
  \added{\Gamma} \mathrel{\added{\vdash}} p \subsumed \patWild
}
\and
\inferrule*[lab={\ruleName{sub-list-list}}]
{
  \added{\Gamma} \mathrel{\added{\vdash}} p_i \subsumed q_i \quad (\forall i)
}
{
  \added{\Gamma} \mathrel{\added{\vdash}} \patList{p_1, \ldots, p_n} \subsumed \patList{q_1, \ldots, q_n}
}
\and
\inferrule*[lab={\ruleName{sub-tuple-tuple}}]
{
  \added{\Gamma} \mathrel{\added{\vdash}} p_i \subsumed q_i \quad (\forall i)
}
{
  \added{\Gamma} \mathrel{\added{\vdash}} \patTuple{p_1, \ldots, p_n} \subsumed \patTuple{q_1, \ldots, q_n}
}
\and
\inferrule*[lab={\ruleName{sub-dict-dict}}]
{
  \set{u_1, \ldots, u_m} \subseteq \set{w_1, \ldots, w_n} \\
  \added{\Gamma} \mathrel{\added{\vdash}} p_i \subsumed q_j \quad (\forall i, j.\ w_i = u_j)
}
{
  \added{\Gamma} \mathrel{\added{\vdash}} \patDict{\bind{w_1}{p_1}, \ldots, \bind{w_n}{p_n}} \subsumed \patDict{\bind{u_1}{q_1}, \ldots, \bind{u_m}{q_m}}
}
\and
\added{\inferrule*[lab={\ruleName{sub-constr}}]
{
  \Gamma \vdash q : C \\
  \Gamma \vdash q' : C' \\
  C' \in \ancestors(C) \\
  \Gamma \vdash \seq{p}[x] \subsumed \seq{p}\,'[x] \quad (\forall x \in \fields(C'))
}
{
  \Gamma \vdash \patConstr{q}{\seq{p}} \subsumed \patConstr{q'}{\seq{p}\,'}
}}
\and
\inferrule*[lab={\ruleName{sub-as-l}}]
{
  \added{\Gamma} \mathrel{\added{\vdash}} p \subsumed q
}
{
  \added{\Gamma} \mathrel{\added{\vdash}} p\;\pyinline{as}\;x \subsumed q
}
\and
\inferrule*[lab={\ruleName{sub-as-r}}]
{
  \added{\Gamma} \mathrel{\added{\vdash}} p \subsumed q
}
{
  \added{\Gamma} \mathrel{\added{\vdash}} p \subsumed q\;\pyinline{as}\;y
}
\end{mathpar}
\caption{Pattern $p$ is subsumed by pattern $p'$ under $\Gamma$}
\label{fig:subsumption}
\end{figure}


\subsection{Patterns}

A pattern list is well-formed (\figref{well-formed-pat}) when each pattern is linear (no repeated variable names) and no earlier pattern subsumes a later one. The \pyinline{match} well-formedness rules use this judgement to enforce reachability and linearity of case patterns.

\begin{figure}
\flushleft \shadebox{$\Gamma \vdash p_1 \cons \ldots \cons p_n$}
\begin{mathpar}
\small
\inferrule*[lab={\ruleName{pat-lit}}]
{ \strut }
{ \Gamma \vdash \ell }
\and
\inferrule*[lab={\ruleName{pat-neg-lit}}]
{ \strut }
{ \Gamma \vdash -n }
\and
\inferrule*[lab={\ruleName{pat-var}}]
{ \strut }
{ \Gamma \vdash x }
\and
\inferrule*[lab={\ruleName{pat-wild}}]
{ \strut }
{ \Gamma \vdash \patWild }
\and
\inferrule*[lab={\ruleName{pat-list}}]
{ \Gamma \vdash p_i \quad (\forall i) }
{ \Gamma \vdash \patList{p_1, \ldots, p_n} }
\and
\inferrule*[lab={\ruleName{pat-tuple}}]
{ \Gamma \vdash p_i \quad (\forall i) }
{ \Gamma \vdash \patTuple{p_1, \ldots, p_n} }
\and
\inferrule*[lab={\ruleName{pat-as}}]
{ \Gamma \vdash p }
{ \Gamma \vdash p\;\pyinline{as}\;x }
\and
\inferrule*[lab={\ruleName{pat-class}}]
{
  \Gamma \vdash q \names d(x_1, \ldots, x_n, x_{n+1}, \ldots, x_{n+m}) \\
  \set{y_1, \ldots, y_m} = \set{x_{n+1}, \ldots, x_{n+m}} \\
  \Gamma \vdash p_i \quad (\forall i)
}
{
  \Gamma \vdash \patClass{q}{p_1, \ldots, p_n, y_1 \mathord{=} p_{n+1}, \ldots, y_m \mathord{=} p_{n+m}}
}
\and
\inferrule*[lab={\ruleName{cases-nil}}]
{ \strut }
{ \Gamma \vdash \varepsilon }
\and
\inferrule*[lab={\ruleName{cases-cons}}]
{
  \Gamma \vdash p_1 \\
  \Gamma \vdash p_2 \cons \ldots \cons p_n \\
  \text{variables in } p_1 \text{ distinct} \\
  p_i \not\subsumed p_1 \quad (\forall i.\, 2 \le i \le n)
}
{
  \Gamma \vdash p_1 \cons p_2 \cons \ldots \cons p_n
}
\end{mathpar}
\caption{Case patterns are well-formed under $\Gamma$}
\label{fig:well-formed-pat}
\end{figure}


\subsection{Blocks and expressions}

Every well-formed statement either \emph{definitely returns}, or \emph{definitely assigns} a context
$\Delta$, which may be empty. The base cases are that \pyinline{return e} definitely returns, and \pyinline{x
= e} definitely assigns $\set{\bind{\pyinline{x}}{\statusDefAssigned}}$. \emph{Result types} $R$ classify these outcomes: $R$ is either $\Returns$ (definitely returns) or $\Assigns{\Delta}$ (definitely assigns $\Delta$). Block well-formedness is given in
\figref{well-formed}; expression well-formedness in \figref{well-formed-expr}.

Context operations $\override$ and $\merge$ lift to a right semiring over result types $R$, with $\Assigns{\varnothing}$ as left and right unit for $\override$ and $\Returns$ as unit for $\merge$ and zero for $\override$.

\begin{definition}[Override and merge on result types]
{\small
\begin{align*}
\Assigns{\Delta} \merge \Assigns{\Delta'} & = \Assigns{(\Delta \merge \Delta')} \\
R \merge \Returns = \Returns \merge R & = R \\
\\
\Assigns{\Delta} \override \Assigns{\Delta'} & = \Assigns{(\Delta \override \Delta')} \\
R \override \Returns = \Returns \override R & = \Returns
\end{align*}}
\end{definition}

\noindent The $\Returns \override R$ case never arises because the well-formedness judgement excludes unreachable statements (\secref{unreachable}).

\begin{lemma}[Right distributivity (result types)]
\[
(R \merge R') \override S = (R \override S) \merge (R' \override S)
\]
\end{lemma}

\noindent A module attribute reference is well-formed only when it resolves to a definitely assigned member
(rule \ruleName{attr-module}); see \secref{no-load-dependent}.

\begin{figure}
\flushleft \shadebox{$\Gamma; \Lambda \vdash b : R$}
\begin{mathpar}
\small
\inferrule*[lab={\ruleName{pass}}]
{
  \strut
}
{
  \Gamma; \Lambda \vdash \exPass : \Assigns{\varnothing}
}
\and
\inferrule*[lab={\ruleName{assign}}]
{
  \Gamma; \Lambda \vdash e \\
  x \notin \capturesE(e)
}
{
  \Gamma; \Lambda \vdash \exAssign{x}{e} : \Assigns{\set{x : \statusDefAssigned}}
}
\and
\inferrule*[lab={\ruleName{expr-stmt}}]
{
  \Gamma; \Lambda \vdash e
}
{
  \Gamma; \Lambda \vdash e : \Assigns{\varnothing}
}
\and
\inferrule*[lab={\ruleName{assert}}]
{
  \Gamma; \Lambda \vdash e
}
{
  \Gamma; \Lambda \vdash \exAssert{e} : \Assigns{\varnothing}
}
\and
\inferrule*[lab={\ruleName{assert-msg}}]
{
  \Gamma; \Lambda \vdash e \\
  \Gamma; \Lambda \vdash e'
}
{
  \Gamma; \Lambda \vdash \exAssertMsg{e}{e'} : \Assigns{\varnothing}
}
\and
\inferrule*[lab={\ruleName{return-none}}]
{
  \strut
}
{
  \Gamma; \Lambda \vdash \exReturnNone : \Returns
}
\and
\inferrule*[lab={\ruleName{return}}]
{
  \Gamma; \Lambda \vdash e
}
{
  \Gamma; \Lambda \vdash \exReturn{e} : \Returns
}
\and
\inferrule*[lab={\ruleName{if}}]
{
  \Gamma; \Lambda \vdash e_i \quad (\forall i) \\
  \Gamma; \Lambda \vdash b_i : R_i \quad (\forall i)
}
{
  \Gamma; \Lambda \vdash \exIf{e_1}{b_1}{\exElif{e_2}{b_2}\;\ldots\;\exElif{e_{n+1}}{b_{n+1}}} : R_1 \merge \ldots R_{n + 1} \merge \ldots \merge \Assigns{\varnothing}
}
\and
\inferrule*[lab={\ruleName{if-else}}]
{
  \Gamma; \Lambda \vdash e_i \quad (\forall i) \\
  \Gamma; \Lambda \vdash b_i : R_i \quad (\forall i)
}
{
  \Gamma; \Lambda \vdash \exIfElse{e_1}{b_1}{\exElif{e_2}{b_2}\;\ldots\;\exElif{e_{n+1}}{b_{n+1}}}{b_{n+2}} : R_1 \merge \ldots \merge R_{n+2}
}
\and
\inferrule*[lab={\ruleName{mutual}}]
{
  \Gamma \runion \statusDefAssigned(\seq{f}) \runion \statusDefAssigned(\seq{x}_{i}) \runion
  \statusNotDefAssigned(\assignsS(b_i) \setminus \seq{x}_i); \Lambda \vdash b_i : R_i
  \quad (\forall i)
  \\
  f_1, \ldots, f_{n+1} \text{ distinct}
}
{
  \Gamma; \Lambda \vdash \exDef{f_1}{\vec{x}_1}{b_1} \cons \ldots \cons \exDef{f_{n+1}}{\vec{x}_{n+1}}{b_{n+1}} : \Assigns{\statusDefAssigned(\seq{f})}
}
\and
\inferrule*[lab={\ruleName{match}}]
{
  \Gamma; \Lambda \vdash e \\
  \Gamma \runion \statusDefAssigned(\binds(p_i)); \Lambda \vdash b_i : R_i \quad (\forall i) \\
  S_i = \Assigns{\statusDefAssigned(\binds(p_i))} \runion R_i \quad (\forall i) \\
  \Lambda \vdash p_1 \cons \ldots \cons p_{n+1} \\
  p_{n+1} \in \set{x, \patWild}
}
{
  \Gamma; \Lambda \vdash \exMatch{e}{\exCase{p_1}{b_1}\;\ldots\;\exCase{p_{n+1}}{b_{n+1}}} : S_1 \merge \ldots \merge S_{n+1}
}
\and
\inferrule*[lab={\ruleName{match-partial}}]
{
  \Gamma; \Lambda \vdash e \\
  \Gamma \runion \statusDefAssigned(\binds(p_i)); \Lambda \vdash b_i : R_i \quad (\forall i) \\
  S_i = \Assigns{\statusDefAssigned(\binds(p_i))} \runion R_i \quad (\forall i) \\
  \Lambda \vdash p_1 \cons \ldots \cons p_{n+1} \\
  p_{n+1} \notin \set{x, \patWild}
}
{
  \Gamma; \Lambda \vdash \exMatch{e}{\exCase{p_1}{b_1}\;\ldots\;\exCase{p_{n+1}}{b_{n+1}}} : S_1 \merge \ldots \merge S_{n+1} \merge \Assigns{\varnothing}
}
\and
\inferrule*[lab={\ruleName{cons}}]
{
  \Gamma; \Lambda \vdash s : \Assigns{\Delta} \\
  \Gamma \runion \Delta; \Lambda \vdash b : R \\
  \capturesS(s) \cap \assignsS(b) = \varnothing
}
{
  \Gamma; \Lambda \vdash s \cons b : \Assigns{\Delta} \runion R
}
\end{mathpar}
\caption{Block $b$ is well-formed in $\Gamma$ and $\Lambda$}
\label{fig:well-formed}
\end{figure}

\begin{figure}
\begin{subfigure}{\textwidth}
\flushleft \shadebox{$\Gamma \vdash q \names \theta$}
\begin{mathpar}
\small
\inferrule*[lab={\ruleName{simple-module}}]
{
  x : q \in \Gamma
}
{
  \Gamma \vdash x \names q
}
\and
\inferrule*[lab={\ruleName{simple-class}}]
{
  c : (q', \vec{y}, c') \in \Gamma \\
  \fields(q'.c) = \vec{x}
}
{
  \Gamma \vdash c \names q'.c(\vec{x})
}
\and
\inferrule*[lab={\ruleName{qualified-module}}]
{
  \Gamma \vdash q \names q' \\
  q'.x \in \dom(\mathcal{M})
}
{
  \Gamma \vdash q.x \names q'.x
}
\and
\inferrule*[lab={\ruleName{qualified-class}}]
{
  \Gamma \vdash q \names q' \\
  \Gamma_{q'}(c) = (q', \vec{y}, c'') \\
  \fields(q'.c) = \vec{x}
}
{
  \Gamma \vdash q.c \names q'.c(\vec{x})
}
\end{mathpar}
\subcaption{Name $q$ resolves to context entry $\theta$ under $\Gamma$}
\end{subfigure}

\vspace{2mm}

\begin{subfigure}{\textwidth}
\flushleft \shadebox{$\Gamma \vdash e$}
\begin{mathpar}
\small
\inferrule*[lab={\ruleName{var}}]
{
  \Gamma(x) = \statusDefAssigned
}
{
  \Gamma \vdash x
}
\and
\inferrule*[lab={\ruleName{lambda}}]
{
  \Gamma \runion \statusDefAssigned(\seq{x}) \vdash e
}
{
  \Gamma \vdash \exLambda{\seq{x}}{e}
}
\and
\inferrule*[lab={\ruleName{constr}}]
{
  \Gamma \vdash q \names d(x_1, \ldots, x_n) \\
  \Gamma \vdash e_i \quad (\forall i)
}
{
  \Gamma \vdash \exConstruct{q}{e_1, \ldots, e_n}
}
\and
\inferrule*[lab={\ruleName{attr-module}}]
{
  \Gamma \vdash e \names q \\
  y \in \dom(\Gamma_q)
}
{
  \Gamma \vdash e.y
}
\and
\inferrule*[lab={\ruleName{attr-object}}]
{
  \Gamma \vdash e
}
{
  \Gamma \vdash e.y
}
\end{mathpar}
\subcaption{Expression $e$ is well-formed under $\Gamma$}
\end{subfigure}
\caption{Expression well-formedness}
\label{fig:well-formed-expr}
\end{figure}


\subsection{Modules and programs}

In this section, loading a module means checking its body and computing its signature
(\figref{well-formed-module}), the static counterpart of run-time loading (\secref{opsem}). Both
import forms load their target, together with its ancestors
(\figref{well-formed-imports}). A plain import also records each of these modules as a member of its
parent (the loads-as judgement), so that the binding created for the root of the target gives attribute
access along the whole chain. A from-import also loads any imported name that denotes a submodule. An
ancestor of the target is exempt from loading first if it is the importing module or one of its
ancestors, since at run time it may still be loading. The exemption applies only to from-imports, because
for a plain import
the only ancestor that could never load first is the importing module itself, and an import naming a proper
descendant of the importing module is excluded outright (\secref{no-load-dependent}). Loading a module
may thus require other modules to load first; the rules are interpreted inductively, so a program with an
import cycle is ill-formed, because a module in a cycle could finish loading only after itself. Program
well-formedness (\figref{well-formed-module}) requires the main module to be well-formed and hence,
recursively, every module the program can load; a module that is never imported is not checked.

\begin{figure}
\flushleft \shadebox{$\Gamma \vdash m : \Delta$}
\begin{mathpar}
\small
\inferrule*[lab={\ruleName{empty}}]
{
  \strut
}
{
  \Gamma \vdash \varepsilon : \varnothing
}
\and
\inferrule*[lab={\ruleName{import}}]
{
  M' = x_1.x_2.\ldots.x_{n+1} \\
  \Gamma \runion \statusDefAssigned(\set{x_1}) \vdash m : \Delta
}
{
  \Gamma \vdash \exImport{M'} \cons m : \statusDefAssigned(\set{x_1}) \runion \Delta
}
\and
\inferrule*[lab={\ruleName{from-import}}]
{
  \Gamma \runion \statusDefAssigned(\set{x_1, \ldots, x_{n+1}}) \vdash m : \Delta
}
{
  \Gamma \vdash \exFromImport{M'}{x_1, \ldots, x_{n+1}} \cons m : \statusDefAssigned(\set{x_1, \ldots, x_{n+1}}) \runion \Delta
}
\and
\inferrule*[lab={\ruleName{statement}}]
{
  \Gamma \vdash s : \Assigns{\Delta} \\
  \Gamma \runion \Delta \vdash m : \Delta' \\
  \capturesS(s) \cap \assignsS(m) = \varnothing
}
{
  \Gamma \vdash s \cons m : \Delta \runion \Delta'
}
\and
\inferrule*[lab={\ruleName{class}}]
{
  \decls(\phi) = x_1, \ldots, x_n \\
  x_1, \ldots, x_n \text{ distinct} \\
  \Gamma \vdash m : \Delta
}
{
  \Gamma \vdash \exDataclass{C}{\phi} \cons m : \Delta
}
\and
\inferrule*[lab={\ruleName{class-extend}}]
{
  B \in \dom(\Lambda_M) \\
  \decls(\phi) = x_1, \ldots, x_n \\
  x_1, \ldots, x_n \text{ distinct and disjoint from } \fields(M.B) \\
  \Gamma \vdash m : \Delta
}
{
  \Gamma \vdash \exDataclassExtend{C}{B}{\phi} \cons m : \Delta
}
\end{mathpar}
\caption{Body $m$ of module $M$ is well-formed in $\Gamma$}
\label{fig:well-formed-module}
\end{figure}


The signature $\Gamma$ a module's load judgement produces (\figref{well-formed-module}) comprises the classes
and variables a module defines, together with its intrinsic \pyinline{__name__} and a stub for each submodule,
never the names it imports: following PEP 484\footnote{\url{https://peps.python.org/pep-0484/}}, an imported
name is a private dependency, not re-exported as part of the interface. Submodule stubs and the module's own
bindings are disjoint, because the module rules exclude a module that binds a name that is also the name of
a submodule (\secref{no-load-dependent}).

\subsection{Examples}

\subsubsection{Assignment as definition}
\label{sec:assignment-as-definition}

Although \PurePy is pure and all variable usage is well-scoped, local variables are introduced using an
\emph{assignment statement}, for compatibility with Python. If a variable is definitely assigned in any branch
of a conditional, then for that variable to be usable after the conditional, every other branch of the
conditional must either definitely return, or definitely assign that variable too. Thus the following are both
well-formed:

\begin{center}
\begin{minipage}{0.48\textwidth}
\centering
\begin{python}
if a < b:
   comp = "smaller"
else:
   comp = "bigger"
return comp
\end{python}
\end{minipage}
\hfill
\begin{minipage}{0.48\textwidth}
\centering
\begin{python}
if a < b:
   return "smaller"
else:
   comp = "bigger"
return comp
\end{python}
\end{minipage}
\end{center}
\medskip

\noindent The definite assignment set for a conditional that has at least one branch which definitely assigns
is calculated by merging the definite assignment sets for all branches which definitely assign, using the
$\merge$ operator on environments.

Definite assignment thus gives each reference a unique lexical referent: a variable is introduced by
an assignment, possibly jointly by the branches of a conditional, and a name may be referenced only where one
such introduction covers every path. This is stricter than avoiding Python's \pyinline{NameError}: the
left-hand program of \secref{contexts} assigns \pyinline{x} on every run, but is rejected because the
reference would denote the branch-local variable on one path and the variable it hides on the other.

\subsubsection{Local assignment in conditional clauses}
\label{sec:local-assignment}

For flexibility, we permit \emph{branch-local} variables. The assignment of \pyinline{x} below (left) is
permitted, and can be used locally in that branch of the conditional:

\begin{center}
\begin{minipage}{0.48\textwidth}
\centering
\begin{python}
def foo():
  z = 6
  if b:
    x = 3
    y = x + 1
  else:
    y = 0
  return z
\end{python}

\smallskip

{\small Branch-local definition of \pyinline{x}}
\end{minipage}
\hfill
\begin{minipage}{0.48\textwidth}
\centering
\begin{python}
def foo():
  x = 6
  if b:
    x = 3
    y = x + 1
  else:
    y = 0
  return x # error: partially defined
\end{python}

\smallskip

{\small Now with use of \pyinline{x} after the conditional}
\end{minipage}
\end{center}
\medskip

\noindent However such variables only appear to be branch-local. The definition of $\merge$ actually ensures
that a variable assigned in any branch is considered assigned by the conditional; however unless it is
assigned in \emph{every} branch that does not definitely return, it will be bound to $\statusNotDefAssigned$ in the
corresponding environment. Any attempt to refer to such a variable is illegal, as it refers to a variable
which may not be assigned. (Thus the right-hand program is ill-formed, despite the enclosing \pyinline{x = 6},
because in the \pyinline{return x} statement, the inner assignment shadows the the outer one, and the inner
assignment is only partial.) If we were to allow the program on the right, there would be no straightforward
way to assign it a semantics which is both pure, and in agreement with Python, since in Python the function
returns \pyinline{3} in the case where \pyinline{b} is \pyinline{True}.

\subsubsection{Function-level scoping of local variables}

In Python, all local variables are scoped to the entire function body, not just from the point of
assignment onwards. The $\assignsS$ function in the \ruleName{def} rule reflects this: all variables
assigned anywhere in the function body are introduced at $\statusNotDefAssigned$ before the body is checked. This means
the following program is statically rejected by \PurePy, whereas Python rejects it at runtime with an
\pyinline{UnboundLocalError}:

\begin{python}
y = 7

def f():
    x = y  # error: y is local (due to assignment below)
    y = 8

f()
\end{python}

\noindent Because \pyinline{y} appears on the left-hand side of an assignment in \pyinline{f}, it is in
$\assignsS(b)$ for the body of \pyinline{f}, and is therefore introduced at $\statusNotDefAssigned$ at the top of the
function. The reference to \pyinline{y} in \pyinline{x = y} then fails the well-formedness check, since
$\Gamma(\pyinline{y}) = \statusNotDefAssigned$ rather than $\Gamma(\pyinline{y}) = \statusDefAssigned$.

\subsubsection{Reassignment of captured variables}
\label{sec:captured-vars}

The exclusion of captured-variable reassignment (\secref{no-mutable-state}) is enforced by a side condition
on the \ruleName{cons} rule: no variable captured by a statement $s$ may be reassigned in the rest of the block
$b$. In the example of \secref{no-mutable-state}, \pyinline{x} is captured by the definition of
\pyinline{g} and reassigned by \pyinline{x = 6}, so the program is rejected. A dataclass declaration counts as
an assignment of its class name for these definitions, though re-declaring a class name is excluded outright
(\secref{no-first-class}).

The same issue arises when a variable is captured and reassigned in the \emph{same} statement:

\begin{python}
x = 5
def x():
    return x
x()  # Python: returns the function (late binding); pure: returns 5
\end{python}

\noindent The side condition on the \ruleName{assign} rule rejects programs like this, requiring that the
variable being bound is not captured by the expression. The \ruleName{def} rule, by contrast, binds
each $f_i$ in the typing context for every body $b_i$, so the function names may legitimately appear in
$\capturesS$ of those bodies. This is what permits self-recursion (the singleton-region case) and
mutual recursion within a region.

\subsubsection{Mutual recursion}
\label{sec:mutual-recursion}

The \ruleName{def} rule binds the names $f_1, \ldots, f_{n+1}$ of a mutual region $d_1 \ldots d_{n+1}$
simultaneously. Each body $b_i$ is checked in a context that includes all the $f_j$ as definitely
assigned, so any $f_j$ may reference any other $f_k$ within the region (the singleton case $n = 0$
permits self-recursion). The names in a region must be distinct.

Inserting any non-\pyinline{def} statement between two \pyinline{def}s splits them into separate regions,
breaking the mutual reference:

\begin{python}
def even(n):
    if n == 0:
        return True
    return odd(n - 1)

x = even(0)  # intervening statement separates the regions

def odd(n):
    if n == 0:
        return False
    return even(n - 1)
\end{python}

\noindent This program is rejected statically: \pyinline{even} forms a singleton region in which
\pyinline{odd} is not yet bound, so the reference to \pyinline{odd} in its body fails the
\ruleName{var} rule. Python accepts it: the call \pyinline{even(0)} does not actually reach
\pyinline{odd}, and Python's late binding means \pyinline{odd} is only looked up when needed.

Each branch of a conditional may contain its own mutual regions, subject to the usual definite assignment
rules. The following is well-formed because both branches define both names:

\begin{python}
b = True

if b:
    def f():
        return g()
    def g():
        return "via mutual region"
else:
    def f():
        return "f independent"
    def g():
        return "g independent"

print(f())
print(g())
\end{python}

A region containing two definitions of the same name is rejected directly by the \ruleName{def}
rule's distinct-names side condition:

\begin{python}
def g():
    return 0

def f():
    return g()
def g():     # error: g defined twice in the same region
    return 1
\end{python}

\noindent (Python accepts this; late binding makes \pyinline{f}'s call to \pyinline{g} resolve to the
second \pyinline{def g}.)

Wrapping the second \pyinline{def f; def g} pair in a conditional sidesteps the duplicate-name issue:
the outer \pyinline{def g} and the inner \pyinline{def g} are now in different blocks, so they no longer
fall in the same region. One might still expect a captured-variable problem, since the inner
\pyinline{def g} could be seen as rebinding the outer \pyinline{g} that \pyinline{f} captures. No such
problem arises:

\begin{python}
def g():
    return 0

if b:
    def f():
        return g()
    def g():
        return 1
\end{python}

\noindent The inner \pyinline{def f; def g} form a branch-local mutual region, so \pyinline{f}'s body
resolves \pyinline{g} to its sibling rather than to the outer \pyinline{g}. The two are bound
simultaneously, so \pyinline{f} cannot be called before the inner \pyinline{g} exists, and the value
\pyinline{f} sees is the same as the value Python's late binding finds at call time.

\subsubsection{Pattern variables escape the match}

A variable bound by a \pyinline{case} pattern is in scope after the \pyinline{match} statement, just as
a variable assigned inside a conditional branch is in scope after the conditional. Python treats matching
this way; \PurePy follows.

\begin{python}
v = 5
match v:
    case x:
        pass
print(x)    # prints 5
\end{python}

Here the catch-all \pyinline{case x:} binds \pyinline{x} in the body's scope, and the binding survives into
the enclosing scope. Definite assignment after the match follows the same merge logic as a conditional: if
every case binds the variable (or returns), it is definitely assigned afterwards; otherwise it is not (see
also \secref{local-assignment}).

\subsubsection{Unreachable code}
\label{sec:unreachable}
A non-empty sequence of statements returns only if the last statement returns, and all preceding statements
assign. Thus the following is illegal:

\begin{python}
def foo(x):
  return x + 1
  return x    # error: unreachable
\end{python}

\begin{DeletedBlock}
\subsubsection{Conditionals without \pyinline{else} clauses}

An \pyinline{if} statement without an \pyinline{else} clause is equivalent to an \pyinline{if} statement with
an \pyinline{else} clause whose body is a $\exPass$ statement; such a statement never returns.

\subsubsection{Implicit \pyinline{None}}

A \pyinline{return} statement without an expression is equivalent to \pyinline{return None}.

\subsubsection{Implicit \pyinline{return}}

A function whose body $s$ assigns rather than returns is equivalent to a function whose body is $s$ followed
by \pyinline{return None}. This ensures every function returns, even if it is not explicitly specified.

\subsubsection{\pyinline{assert} without a message}

An \pyinline{assert e} statement without a message is equivalent to \pyinline{assert e, ""}.
\end{DeletedBlock}

\section{Operational semantics}
\label{sec:opsem}

\subsection{Values and environments}

Values, environments, and evaluation results are given in \figref{values}. For a statement $s$, the result
$\assigns{\rho}$ represents \emph{normal completion} of $s$, with accumulated environment $\rho$; the result
$\returns{v}$ represents \emph{abrupt completion} of $s$, with return value $v$; \Replaced{and the result
$\aborts{\kappa}$ represents abrupt completion of the whole run, with termination kind $\kappa$}{and
$\aborts{\kappa}$ abrupt completion with an error}.
Composing statements sequentially composes their results according to the following partial monoid, analogous
to sequential composition of result types:
{\small
\begin{align*}
\assigns{\rho} \override \assigns{\rho'} & = \assigns{(\rho \override \rho')} \\
\assigns{\rho} \override \returns{v} & = \returns{v} \\
\assigns{\rho} \override \aborts{\kappa} & = \aborts{\kappa}
\end{align*}}

with $\assigns{\varnothing}$ as left and right unit.

\begin{figure}
\small
\begin{center}
\begin{tabular}{rcll}
\grammarGroup{Values} \\
$v$ & $::=$ & $\ell$ & (Literal) \\
    & $\mid$ & $\exList{v_1, \ldots, v_n}$ & (List) \\
    & $\mid$ & $\exTuple{v_1, \ldots, v_n}$ & (Tuple) \\
    & $\mid$ & $\exDict{\delta}$ & (Dict) \\
    & $\mid$ & $\exLamClosure[\Gamma]{\rho}{\seq{x}}{e}$ & (LambdaClosure) \\
    & $\mid$ & $\exDefClosure[\Gamma]{\rho}{\seq{d}}{i}$ & (DefClosure) \\
    & $\mid$ & $\exObj{\Gamma}{C}{\rho}$ & (Object) \\
    & $\mid$ & $\modStub{q}$ & (ModuleStub) \\
    & $\mid$ & $\modLoaded{q}{\rho}$ & (LoadedModule) \\
    & $\mid$ & $\clsEntry{\Gamma}{q.c}{\seq{x}}{c'}$ & (Class) \\
    & $\mid$ & $\primVal{\primFn}$ & (Primitive)
   \\[3mm]
\grammarGroup{Environments} \\
$\rho$ & $::=$ & $\set{\bind{x_1}{v_1}, \ldots, \bind{x_n}{v_n}}$ & 
   \\[3mm]
\grammarGroup{Dictionary entries} \\
$\delta$ & $::=$ & $\bind{w_1}{v_1} \cons \ldots \cons \bind{w_n}{v_n}$ &
   \\[3mm]
\grammarGroup{Results} \\
$r$ & $::=$ & $\assigns{\rho}$ & (Assigns) \\
  & $\mid$ & $\returns{v}$ & (Returns) \\
  & $\mid$ & $\aborts{\kappa}$ & (Aborts)
   \\[3mm]
\grammarGroup{Outcomes} \\
$o$ & $::=$ & $\val{v}$ & (Val) \\
  & $\mid$ & $\aborts{\kappa}$ & (Aborts)
   \\[3mm]
\grammarGroup{Match results} \\
$m$ & $::=$ & $\matchOk{\rho}$ & (Match) \\
  & $\mid$ & $\noMatch$ & (NoMatch)
   \\[3mm]
\grammarGroup{Termination kinds} \\
$\kappa$ & $::=$ & $\errType \mid \errIndex \mid \errKey \mid \errZeroDiv$ &  \\
    & $\mid$ & $\errAssertNoMsg \mid \errAssert{w} \mid \errAttr \mid \errExit{n}$ &
\end{tabular}
\end{center}
\caption{Values, environments and evaluation results}
\label{fig:values}
\end{figure}


\Added{The primitive $\primVal{\primFn}$ is how $\mathcal{P}$ supplies a member that can be called, since a
supplied member has no body to evaluate and $\primFn$ gives its meaning directly. No rule applies a
primitive, and \ruleName{eval-call-nonfun} does not cover one, so calling a primitive has no derivation
(\issue{132}, \issue{47}).}

A value is well-formed at a context entry (\figref{well-formed-env}): ordinary values at
$\statusDefAssigned$, module references and class entries at the corresponding entry. An environment is
well-formed at a context that binds its variables at corresponding entries.

\begin{figure}
\flushleft \shadebox{$\vdash v$}
\begin{mathpar}
\small
\inferrule*[lab={\ruleName{val-lit}}]
{
  \strut
}
{
  \vdash \ell
}
\and
\inferrule*[lab={\ruleName{val-list}}]
{
  \vdash v_i \quad (\forall i)
}
{
  \vdash \exList{v_1, \ldots, v_n}
}
\and
\inferrule*[lab={\ruleName{val-tuple}}]
{
  \vdash v_i \quad (\forall i)
}
{
  \vdash \exTuple{v_1, \ldots, v_n}
}
\and
\inferrule*[lab={\ruleName{val-lambda}}]
{
  \Gamma \vdash \rho \\
  \Gamma \runion \statusDefAssigned(\seq{x}) \vdash e
}
{
  \vdash \exClosure[\Gamma]{\rho}{\seq{x}}{e}
}
\and
\inferrule*[lab={\ruleName{val-def}}]
{
  \Gamma \vdash \rho \\
  d_j = \exDef{f_j}{\vec{x}_j}{b_j} \quad (\forall j) \\
  \Gamma \runion \statusDefAssigned(\seq{f}) \runion \statusDefAssigned(\seq{x}_j) \runion
  \statusNotDefAssigned(\assignsS(b_j) \setminus \seq{x}_j) \vdash b_j : R_j \quad (\forall j) \\
  \seq{f} \text{ distinct}
}
{
  \vdash \exDefClosure[\Gamma]{\rho}{\seq{d}}{i}
}
\and
\inferrule*[lab={\ruleName{val-obj}}]
{
  \Gamma \vdash \rho
}
{
  \vdash \exObj[\Gamma]{d}{\rho}
}
\end{mathpar}

\vspace{2mm}

\flushleft \shadebox{$\Gamma \vdash \rho$}
\begin{mathpar}
\small
\inferrule*[lab={\ruleName{env-var}}]
{
  x : a \in \Gamma \\
  \vdash v
}
{
  \Gamma \vdash \set{x \mapsto v}
}
\and
\inferrule*[lab={\ruleName{env-module}}]
{
  x : \modEntry{q} \in \Gamma \\
  \Delta \vdash \rho'
}
{
  \Gamma \vdash \set{x \mapsto \exObj[\Delta]{\bot}{\rho'}}
}
\and
\inferrule*[lab={\ruleName{env-class}}]
{
  x : \clsEntry{q}{\vec{y}}{c} \in \Gamma
}
{
  \Gamma \vdash \set{x \mapsto \clsEntry{q}{\vec{y}}{c}}
}
\and
\inferrule*[lab={\ruleName{env-cons}}]
{
  \Gamma \vdash \set{x \mapsto w} \\
  \Gamma \vdash \rho
}
{
  \Gamma \vdash \set{x \mapsto w} \runion \rho
}
\end{mathpar}
\caption{Well-formed values and environments}
\label{fig:well-formed-env}
\end{figure}


The rules are given in a monadic metalanguage over outcomes (\figref{values}), in which
$\mlet{v}{o}{\ldots}$ binds the value of an outcome; for example, \ruleName{eval-binop} concludes
$\mlet{v'}{o}{\hat{\odot}(v, v')}$, which is $\hat{\odot}(v, v')$ where $o$ is $\val{v'}$, and
$\aborts{\kappa}$ where $o$ is $\aborts{\kappa}$. A final
premise needs no propagation rule; earlier premises have one each, recording which subexpressions succeeded
before the failing one was reached (\figref{eval-fails}). The statement, sequence and module judgements use
outcomes of the same form, carrying a result, a sequence of values or an environment in place of a value.

\subsection{Auxiliary definitions}

The strict operators are given by metafunctions $\hat{\odot}$ and $\hat{\ominus}$ (\figref{operators}), in
terms of structural equality $\eqOp$ and membership $\memOp$ (\figrefTwo{equality}{membership}). They have
the meanings they have in Python and yield $\aborts{\kappa}$ where Python raises. All three are undefined
where no case applies, and where one of them is undefined the expression has no derivation, which is where
\PurePy declines to give the operation the meaning Python gives it; a non-Boolean test is the same case
(\secref{no-coercions}, \issue{162}).

\begin{figure}
\small
\begin{center}
\begin{tabular}{>{\raggedright\arraybackslash}p{0.22\textwidth}>{\raggedright\arraybackslash}p{0.62\textwidth}}
\hline
$\odot$ & $\hat{\odot}(v, v')$ \\
\hline
\pyinline{==} & $\val{\pyinline{False}}$ if both operands are \pyinline{nan}, and $\val{\eqOp(v, v')}$ otherwise \\
\pyinline{!=} & $\val{\pyinline{True}}$ if both operands are \pyinline{nan} or $\eqOp(v, v') = \pyinline{False}$, and $\val{\pyinline{False}}$ if $\eqOp(v, v') = \pyinline{True}$ \\
\pyinline{in} & $\val{\memOp(v', v)}$ \\
\pyinline{not in} & $\val{\pyinline{True}}$ if $\memOp(v', v) = \pyinline{False}$, and $\val{\pyinline{False}}$ if $\memOp(v', v) = \pyinline{True}$ \\
\pyinline{<}\quad\pyinline{<=}\quad\pyinline{>}\quad\pyinline{>=} & \todo{Ordering, as in Python} \\
\pyinline{+}\quad\pyinline{-}\quad\pyinline{*}\quad\pyinline{/}\quad\pyinline{//}\quad\pyinline{\%}\quad\pyinline{**} & \todo{Arithmetic, as in Python. Float precision is left to the implementation. $\aborts{\errZeroDiv}$ where the divisor is zero, and $\aborts{\errType}$ where an operand is not a number} \\
\hline
\end{tabular}


\vspace{2mm}

\input{functions/unop}
\end{center}
\caption{Operator meanings. An operator returns an outcome}
\label{fig:operators}
\end{figure}

\begin{figure}
\small
\begin{align*}
\eqOp(\ell, \ell') &{}= \pyinline{True} && \text{if } \ell = \ell' \\
\eqOp(\ell, \ell') &{}= \pyinline{False} && \text{if } \ell \neq \ell' \text{ and } \ell, \ell' \text{ have the same kind} \\
\\
\eqOp(\exList{v_1, \ldots, v_n}, \exList{v_1', \ldots, v_n'}) &{}= \textstyle\bigwedge_{1 \leq i \leq n} \eqOp(v_i, v_i') \\
\eqOp(\exTuple{v_1, \ldots, v_n}, \exTuple{v_1', \ldots, v_n'}) &{}= \textstyle\bigwedge_{1 \leq i \leq n} \eqOp(v_i, v_i') \\
\eqOp(\exList{\seq{v}}, \exList{\seq{v}\,'}) = \eqOp(\exTuple{\seq{v}}, \exTuple{\seq{v}\,'}) &{}= \pyinline{False} && \text{if } \seq{v}, \seq{v}\,' \text{ differ in length} \\
\\
\eqOp(d, d') &{}= \textstyle\bigwedge_{w \in \dom(d)} \eqOp(d(w), d'(w)) && \text{if } \dom(d) = \dom(d') \\
\eqOp(d, d') &{}= \pyinline{False} && \text{if } \dom(d) \neq \dom(d') \\
\\
\eqOp(\exObj{\Gamma}{C}{\rho}, \exObj{\Gamma'}{C}{\rho'}) &{}= \textstyle\bigwedge_{x \in \fields(C)} \eqOp(\rho(x), \rho'(x)) \\
\eqOp(\exObj{\Gamma}{C}{\rho}, \exObj{\Gamma'}{C'}{\rho'}) &{}= \pyinline{False} && \text{if } C \neq C'
\end{align*}
\caption{Structural equality; undefined on operands of differing kinds and on closures}
\label{fig:equality}
\end{figure}

\begin{figure}
\small
\begin{align*}
\memOp(\exList{v_1, \ldots, v_n}, v) &{}= \memElems(v_1 \cons \ldots \cons v_n, v) \\*
\memOp(\exTuple{v_1, \ldots, v_n}, v) &{}= \memElems(v_1 \cons \ldots \cons v_n, v) \\*
\\
\memOp(\delta, w) &{}= \pyinline{True} && \text{if } w \in \dom(\delta) \\*
\memOp(\delta, w) &{}= \pyinline{False} && \text{if } w \notin \dom(\delta) \\*
\\
\memOp(w', w) &{}= \pyinline{True} && \text{if } w' = w_1 \concat w \concat w_2 \text{ for some } w_1, w_2 \\*
\memOp(w', w) &{}= \pyinline{False} && (\text{otherwise})
 \\
\\
\memElems(\varepsilon, v) &{}= \pyinline{False} \\*
\memElems(v' \cons \seq{v}, v) &{}= \pyinline{True} && \text{if } \eqOp(v, v') = \pyinline{True} \\*
\memElems(v' \cons \seq{v}, v) &{}= \memElems(\seq{v}, v) && \text{if } \eqOp(v, v') = \pyinline{False}

\end{align*}
\caption{Membership, undefined where no case applies}
\label{fig:membership}
\end{figure}


\Added{The remaining metafunctions the rules use are collected in \figref{eval-aux}: $\outcome$ reads a
block's result as the outcome of calling it, $\iterOp$ gives the sequence a generator draws from, $\getitem$
interprets subscripting, $\dispatch$ selects the case a $\pyinline{match}$ takes, and $\entries$ and
$\update$ build a dictionary from its entries.}

\begin{figure}
\small
\begin{align*}
\elems(\exList{v_1, \ldots, v_n}) &{}= v_1 \ldots v_n \\
\elems(\exTuple{v_1, \ldots, v_n}) &{}= v_1 \ldots v_n \\
\\
\added{\iterOp(v)} &{}\added{{}= \elems(v)} \\
\\
\added{\getitem(v, n)} &{}\added{{}= \val{v_{n+1}}} && \added{\text{if } \elems(v) = v_1 \ldots v_m \text{ and } 0 \leq n < m} \\
\added{\getitem(v, n)} &{}\added{{}= \val{v_{m+n+1}}} && \added{\text{if } \elems(v) = v_1 \ldots v_m \text{ and } {-m} \leq n < 0} \\
\added{\getitem(v, n)} &{}\added{{}= \fails{\errIndex}} && \added{\text{if } \elems(v) = v_1 \ldots v_m \text{ and } (n < -m \text{ or } n \geq m)} \\
\added{\getitem(v, v')} &{}\added{{}= \fails{\errType}} && \added{(\text{otherwise})} \\
\\
\outcome(\returns{v}) &= \val{v} \\
\outcome(\assigns{\rho}) &= \val{\pyinline{None}} \\
\outcome(\fails{\kappa}) &= \fails{\kappa} \\
\\
\env(\rho, x) &= \rho' && \text{where } \rho(x) = \modLoaded{q'}{\rho'} \\
\env(\rho, q.x) &= \rho' && \text{where } \env(\rho, q)(x) = \modLoaded{q'}{\rho'} \\
\\
\resolveClass(\rho, c) &= \rho(c) \\
\resolveClass(\rho, q.c) &= \env(\rho, q)(c) \\
\\
\dispatch(\rho, v, \varepsilon, \varepsilon) &= (\varnothing, \exPass) \\
\dispatch(\rho, v, p \cons \seq{p}, b \cons \seq{b}) &= (\rho', b) && \text{if } \rho, p, v \matches \rho' \\
\dispatch(\rho, v, p \cons \seq{p}, b \cons \seq{b}) &= \dispatch(\rho, v, \seq{p}, \seq{b}) && \text{if } \rho, p, v \notmatches \\
\\
\entries(\delta, \varepsilon) &{}= \delta \\
\entries(\delta, \bind{w}{v} \cons \delta') &{}= \entries(\update(\delta, w, v), \delta') \\
\\
\update(\varepsilon, w, v) &{}= \bind{w}{v} \\
\update(\bind{w}{u} \cons \delta, w, v) &{}= \bind{w}{v} \cons \delta \\
\update(\bind{w'}{u} \cons \delta, w, v) &{}= \bind{w'}{u} \cons \update(\delta, w, v) && \text{if } w \neq w'
\end{align*}
\caption{Metafunctions used by block and expression evaluation}
\label{fig:eval-aux}
\end{figure}


\subsection{Pattern matching}

The pattern-matching judgement $\rho, p, v \matches m$ (\figref{match}) is a deterministic partial
function. \Added{A \emph{match result} $m$ is either $\matchOk{\rho'}$, a match binding $\rho'$, or
$\noMatch$, a failure to match, on which $\dispatch$ moves to the next case. Having no derivation is a third
outcome, distinct from both: the match has no meaning, and \PurePy gives the run no result. That is where a
pattern and a value are not comparable, so \PurePy can say neither that they match nor that they do not; a
list pattern against a tuple is such a case (\issue{151}). Sub-results compose with $\override$, matching
sequentially:}
{\small
\begin{align*}
\added{\matchOk{\rho} \override \matchOk{\rho'}} & \added{{}= \matchOk{(\rho \disjunion \rho')}} \\
\added{m \override \noMatch = \noMatch \override m} & \added{{}= \noMatch}
\end{align*}}

\Added{with $\matchOk{\varnothing}$ as left and right unit.}
A constructor pattern matches any object whose class has the
pattern's class among its ancestors; sub-patterns bind the fields of the pattern's class. An object records
its class entry, whose qualified name is the identity of the class: since a module declares each class name at
most once, the entry comparison in \ruleName{eval-pat-constr} amounts to a comparison of qualified names.

\begin{figure}
\flushleft \shadebox{$\rho \vdash p, v \matches \rho'$}
\begin{mathpar}
\small
\inferrule*[lab={\ruleName{pat-lit}}]
{
  \strut
}
{
  \rho \vdash \ell, \ell \matches \varnothing
}
\and
\inferrule*[lab={\ruleName{pat-neg-lit}}]
{
  \strut
}
{
  \rho \vdash -n, -n \matches \varnothing
}
\and
\inferrule*[lab={\ruleName{pat-var}}]
{
  \strut
}
{
  \rho \vdash x, v \matches \set{x \mapsto v}
}
\and
\inferrule*[lab={\ruleName{pat-wild}}]
{
  \strut
}
{
  \rho \vdash \patWild, v \matches \varnothing
}
\and
\inferrule*[lab={\ruleName{pat-list}}]
{
  \rho \vdash p_i, v_i \matches \rho_i \quad (\forall i)
}
{
  \rho \vdash \patList{p_1, \ldots, p_n}, \exList{v_1, \ldots, v_n} \matches \rho_1 \disjunion \ldots \disjunion \rho_n
}
\and
\inferrule*[lab={\ruleName{pat-tuple}}]
{
  \rho \vdash p_i, v_i \matches \rho_i \quad (\forall i)
}
{
  \rho \vdash \patTuple{p_1, \ldots, p_n}, \exTuple{v_1, \ldots, v_n} \matches \rho_1 \disjunion \ldots \disjunion \rho_n
}
\and
\inferrule*[lab={\ruleName{pat-as}}]
{
  \rho \vdash p, v \matches \rho'
}
{
  \rho \vdash p\;\pyinline{as}\;x, v \matches \rho' \disjunion \set{x \mapsto v}
}
\and
\inferrule*[lab={\ruleName{pat-class}}]
{
  \fields(\Gamma_q(c)) = x_1, \ldots, x_n, x_{n+1}, \ldots, x_{n+m} \\
  \rho \vdash p_i, \rho_o(x_i) \matches \rho_i \quad (\forall i.\, 1 \le i \le n) \\
  \rho \vdash p_{n+j}, \rho_o(y_j) \matches \rho_{n+j} \quad (\forall j.\, 1 \le j \le m)
}
{
  \rho \vdash \patClass{c}{p_1, \ldots, p_n, y_1 \mathord{=} p_{n+1}, \ldots, y_m \mathord{=} p_{n+m}}, \exObj{q.c}{\rho_o} \matches \rho_1 \disjunion \ldots \disjunion \rho_{n+m}
}
\end{mathpar}
\caption{Pattern matching}
\label{fig:match}
\end{figure}


\subsection{Blocks and expressions}

Block evaluation $\rho, b \Rightarrow r$\Added{, for any $\rho : \Gamma$ and any $b$ well-formed in $\Gamma$,}
is given in \figref{eval-block-outcome}. Function definitions delimit
the scope of sequential-composition effects, so as black boxes functions are always pure. The metafunction $\outcome$ (\figref{eval-aux}) reads a block's result as
the outcome of calling it, $\val{\pyinline{None}}$ where the block falls off the end.

Expression evaluation $\rho, e \Rightarrow o$ (\figrefTwo{eval-outcome}{eval-outcome-ctd})\Added{, for any $\rho : \Gamma$ and any $e$ well-formed in $\Gamma$,} yields an
outcome, either $\val{v}$ or $\aborts{\kappa}$ (\figref{values}). \Replaced{The termination kinds other than $\errExit{n}$ are named after the
exceptions the same program raises under Python in the same circumstances.}{The error kinds are designed to
correspond to the error the same program would raise under Python in the same circumstances.} The operators
\pyinline{and} and \pyinline{or} evaluate their second operand only when the first leaves the result open, and
a conditional expression requires its condition to be \pyinline{True} or \pyinline{False}, as
\ruleName{eval-if} does. A comprehension evaluates its body once for each binding its qualifiers produce, and
concatenates the results; a generator draws from $\iterOp$ (\figref{eval-aux}), which is
\Replaced{defined on lists, tuples, strings and dictionaries, a dictionary yielding its keys as in Python.
These are the only values a generator can draw from.}{defined on lists and tuples only, since \PurePy has no
wider notion of iteration.}

\begin{figure}
\small
\begin{mathpar}
\inferrule*[lab={\ruleName{eval-cond-true}}]
{
  \rho, e \Rightarrow \val{\pyinline{True}} \\
  \rho, e_1 \Rightarrow o
}
{
  \rho, \exCondExpr{e_1}{e}{e_2} \Rightarrow o
}
\end{mathpar}
\caption{\todo{Expression $e$ evaluates to outcome $o$ under environment $\rho : \Gamma$, where $o$ is
$\val{v}$ or $\fails$ (\issue{63})}}
\label{fig:eval-outcome}
\end{figure}


\subsection{Modules and programs}

Module and program evaluation are given in \figref{eval-modules}; import evaluation in \figref{eval-import-prefix}; importing names from a loaded module in
\figref{eval-imports}. Extend (Definition~\ref{def:extend}) extends to module values by reading its equations with a loaded reference $\modLoaded{q}{\rho}$ carrying an environment in place of a context. The loading judgement $q \Rightarrow_{\mathsf{load}} o$ evaluates the import prefix
of $q$ and then its statements, mirroring the well-formedness judgement $\vdash_{\mathsf{load}} q : \Gamma$
(\figref{well-formed-module}); as with the static rules, the mutual
recursion between loading and evaluation is grounded by program well-formedness, which rules out
import cycles (\secref{modules-programs}).
Loading takes a single module; the ancestors of an imported module are loaded by the
loads-as judgement for an import, and directly as premises of \ruleName{eval-from-import} for a
from-import.

\Deleted{The correspondence is exact: if $\vdash_{\mathsf{load}} q : \Gamma$ and $q \Rightarrow_{\mathsf{load}} \val{\rho}$ then $\vdash \rho : \Gamma$ (\figref{well-formed-env}).} \begin{theorem}[\Added{Load preservation}]
\Added{If $\vdash_{\mathsf{load}} q : \Gamma$ and $q \Rightarrow_{\mathsf{load}} \val{\rho}$ then $\vdash \rho : \Gamma$.}
\end{theorem}

\noindent Since
imports precede statements, every module a statement can reference is loaded before the statement runs, and
\Replaced{loading is deterministic, so}{since loading has no observable effects, reloading a module is
unobservable; the semantics therefore needs no cache of loaded modules.} \Added{the semantics needs no cache
of loaded modules. The rule \ruleName{eval-module} puts $\mathcal{P}(\pyinline{builtins})$ under a module
body's environment, mirroring \ruleName{module} (\secref{wf-modules}).}\Deleted{These rules use the environment $\mathcal{P}(\pyinline{builtins})$ binding the names of
\figref{predefined}.}
\Replaced{A program evaluates to a \emph{termination kind} (\figref{values}): $\errExit{0}$ if the main module
loads, and otherwise the kind it aborted with. $\errExit{n}$ is a run that exited with status $n$; the remaining kinds are named after the exceptions the
same program raises under Python. An implementation must agree on which kind a run yields, though how it
reports one is not prescribed.}{A program evaluates to
its exit status: a run whose main module loads successfully yields $0$, and one whose main module fails yields
$1$, as Python does for an uncaught exception.}

\begin{figure}
\begin{subfigure}{\textwidth}
\flushleft \shadebox{$\mu, \rho, m \Rightarrow \mu', \rho'$}
\begin{mathpar}
\small
\inferrule*[lab={\ruleName{eval-nil}}]
{
  \strut
}
{
  \mu, \rho, \varepsilon \Rightarrow \mu, \varnothing
}
\and
\inferrule*[lab={\ruleName{eval-import}}]
{
  \mu, q \Rightarrow_{\mathsf{load}} \mu' \\
  q = x_1.x_2.\ldots.x_{n+1} \\
  \mu', \rho \runion \set{\bind{x_1}{\exMod{\Gamma_{x_1}}{\mu'(x_1)}}}, m \Rightarrow \mu'', \rho'
}
{
  \mu, \rho, \exImport{q} \cons m \Rightarrow \mu'', \rho'
}
\and
\inferrule*[lab={\ruleName{eval-from-import}}]
{
  \mu, q \Rightarrow_{\mathsf{load}} \mu' \\
  \mu', q, \vec{x} \importsR \mu'', \rho' \\
  \mu'', \rho \runion \rho', m \Rightarrow \mu''', \rho''
}
{
  \mu, \rho, \exFromImport{q}{\vec{x}} \cons m \Rightarrow \mu''', \rho''
}
\and
\inferrule*[lab={\ruleName{eval-statement}}]
{
  \rho, s \Rightarrow \assigns{\rho'} \\
  \mu, \rho \runion \rho', m \Rightarrow \mu', \rho''
}
{
  \mu, \rho, s \cons m \Rightarrow \mu', \rho' \runion \rho''
}
\and
\inferrule*[lab={\ruleName{eval-class}}]
{
  C = \clsEntry{\Gamma}{q}{\ownFields(\phi)}{\bot} \\
  \mu, \rho \runion \set{\bind{c}{C}}, m \Rightarrow \mu', \rho'
}
{
  \mu, \rho, \exDataclass{c}{\phi} \cons m \Rightarrow \mu', \set{\bind{c}{C}} \runion \rho'
}
\and
\inferrule*[lab={\ruleName{eval-class-extend}}]
{
  C = \clsEntry{\Gamma}{q}{\ownFields(\phi)}{c'} \\
  \mu, \rho \runion \set{\bind{c}{C}}, m \Rightarrow \mu', \rho'
}
{
  \mu, \rho, \exDataclassExtend{c}{c'}{\phi} \cons m \Rightarrow \mu', \set{\bind{c}{C}} \runion \rho'
}
\end{mathpar}
\subcaption{Module body $m$ of module $q$ evaluates under store $\mu$ and environment $\rho : \Gamma$}
\end{subfigure}

\vspace{2mm}

\begin{subfigure}{\textwidth}
\flushleft \shadebox{$\mathcal{M} \Rightarrow n$}
\begin{mathpar}
\small
\inferrule*[lab={\ruleName{eval-program}}]
{
  \varnothing, \varnothing, (\exAssign{\pymath{\_\_name\_\_}}{\pymath{"\_\_main\_\_"}}) \cons \mathcal{M}(\pymath{\_\_main\_\_}) \Rightarrow \mu, \rho
}
{
  \mathcal{M} \Rightarrow 0
}
\end{mathpar}
\subcaption{Program $\mathcal{M}$ evaluates to exit code $n$}
\end{subfigure}
\caption{Module and program evaluation}
\label{fig:eval-modules}
\end{figure}

\begin{figure}
\begin{subfigure}{\textwidth}
\flushleft \shadebox{$\mu, q \Rightarrow_{\mathsf{load}} \mu'$}
\begin{mathpar}
\small
\inferrule*[lab={\ruleName{load-cached}}]
{
  q \in \dom(\mu)
}
{
  \mu, q \Rightarrow_{\mathsf{load}} \mu
}
\and
\inferrule*[lab={\ruleName{load-unqualified}}]
{
  x \notin \dom(\mu) \\
  \mu, \varnothing, (\exAssign{\pymath{\_\_name\_\_}}{\nameOf(x)}) \cons \mathcal{M}(x) \Rightarrow \mu', \rho
}
{
  \mu, x \Rightarrow_{\mathsf{load}} \mu' \runion \set{\bind{x}{\rho}}
}
\and
\inferrule*[lab={\ruleName{load-qualified}}]
{
  q.x \notin \dom(\mu) \\
  \mu, q \Rightarrow_{\mathsf{load}} \mu' \\
  \mu', \varnothing, (\exAssign{\pymath{\_\_name\_\_}}{\nameOf(q.x)}) \cons \mathcal{M}(q.x) \Rightarrow \mu'', \rho
}
{
  \mu, q.x \Rightarrow_{\mathsf{load}} \mu'' \runion \set{\bind{q.x}{\rho}}
}
\end{mathpar}
\subcaption{Module $q$ is loaded into the store $\mu$}
\end{subfigure}

\vspace{2mm}

\begin{subfigure}{\textwidth}
\flushleft \shadebox{$\mu, q, \vec{x} \importsR \mu', \rho$}
\begin{mathpar}
\small
\inferrule*[lab={\ruleName{imp-nil}}]
{
  \strut
}
{
  \mu, q, \varepsilon \importsR \mu, \varnothing
}
\and
\inferrule*[lab={\ruleName{imp-var}}]
{
  x \in \dom(\mu(q)) \\
  \mu, q, \vec{x} \importsR \mu', \rho
}
{
  \mu, q, x \cons \vec{x} \importsR \mu', \set{\bind{x}{\mu(q)(x)}} \runion \rho
}
\and
\inferrule*[lab={\ruleName{imp-class}}]
{
  x \notin \dom(\mu(q)) \\
  q.x \notin \dom(\mathcal{M}) \\
  \mu, q, \vec{x} \importsR \mu', \rho
}
{
  \mu, q, x \cons \vec{x} \importsR \mu', \rho
}
\and
\inferrule*[lab={\ruleName{imp-module}}]
{
  x \notin \dom(\mu(q)) \\
  q.x \in \dom(\mathcal{M}) \\
  \mu, q.x \Rightarrow_{\mathsf{load}} \mu' \\
  \mu', q, \vec{x} \importsR \mu'', \rho
}
{
  \mu, q, x \cons \vec{x} \importsR \mu'', \set{\bind{x}{\mu'(q.x)}} \runion \rho
}
\end{mathpar}
\subcaption{Names $\vec{x}$ imported from module $q$ produce bindings $\rho$}
\end{subfigure}
\caption{Module loading and importing}
\label{fig:eval-load}
\end{figure}


\section{Reference implementation}
\label{sec:reference-impl}

A reference checker, written in Python over the \pyinline{ast} module, decides membership of \PurePy and
reports why a program is rejected. This section describes the checker and its correspondence with the
well-formedness rules of Section~\ref{sec:syntax}.

\subsection{Testing strategy}
\label{sec:testing}

A conformance test suite pairs each test with the verdict Python itself gives, and runs on CPython, PyPy,
GraalPy, and Pyodide. This section sets out the test taxonomy and the differential and metamorphic strategies
used to build it.

\section{Extensions}
\label{sec:extensions}

\PurePy admits conservative extensions that add expressiveness without compromising purity. This section
presents partial application and piecewise function definitions.

\section{Case studies}
\label{sec:case-studies}

This section reports two case studies that use \PurePy as a common functional language.

\section{Evaluation}
\label{sec:evaluation}

This section evaluates \PurePy empirically: how readily large language models produce programs within the
subset, and how readily existing Python converts to it.

\section{Related work}

Subsets of programming languages have been defined for various purposes in the research literature:

Even though the syntax of programming languages if often described precisely, essentially none of the
languages in mainstream use today formally define their behaviour. Therefore, various well-defined 
subsets of programming languages were constructed to equip languages with formal (often even mechanised) semantics. 
Examples of this include work on Java~\cite{igarashi_featherweight_2001}, OCaml~\cite{seassau_formal_2025}, 
JavaScript~\cite{bodin_trusted_2014,guha_essence_2010,park_kjs_2015,politz_tested_2012}, 
PHP~\cite{hutchison_executable_2014}, Rust~\cite{jung_stacked_2020,jung_rustbelt_2018}, and even lower-level 
languages and intermediate representations such as WebAssembly~\cite{watt_mechanising_2021,huisman_two_2021,watt_wasmref-isabelle_2023} 
and LLVM-IR~\cite{zakowski_modular_2021,zhao_formalizing_2012}. A particular success story of this line of work is the  formalisation of a 
substantial subset of the C programming language, called Clight~\cite{blazy_mechanized_2009}, which lead to the development of a
verified compiler~\cite{leroy_formal_2009}. In a similar manner, the CakeML project~\cite{kumar_cakeml_2014} produced a verified
implementation of a fragment of the ML language~\cite{milner_definition_1997}, based on previous efforts~\cite{lee_towards_2007} of 
mechanising its semantics. Previous work on formalising the semantics of Python include an embedding~\cite{guth_formal_2013} within the 
$\mathbb{K}$ semantics framework~\cite{rosu_overview_2010}, a definition of structural operational semantics~\cite{kohl_executable_2021}, 
and a reduction of most of the language to a core calculus~\cite{politz_python_2013}.

A different line of work concerns the delineation of language subsets for a particular purpose. A prominent example is the 
\emph{Source} family of languages~\cite{anderson_shrinking_2021}, which are progressively extended subsets of JavaScript specifically
geared towards introducing first-year computer science students to the foundations of computing. It is based on a new edition of a
classic introductory text book~\cite{abelson_structure_2022}, in which all program examples (which had previously been in Scheme)
are illustrated in JavaScript. By successively increasing the allowed language fragment with each chapter (corresponding to different
versions of the Source language), it allows to put focus on the particular aspects the authors are trying to convey, while disallowing
students with prior knowledge of programming to take short-cuts by utilising more complex language features. The authors 
of~\cite{anderson_shrinking_2021} also remark, that the choice of JavaScript as a language was a pragmatic one: 
A previous study~\cite{simon_language_2018} showed that \emph{industrial relevance} is perceived as the second-most important factor 
in choosing a programming language for introductory courses by educators in Australiasia and the UK.

Likewise, various subsets of Python have been defined for different purposes in the past. ChocoPy~\cite{padhye_chocopy_2019} describes a
subset specifically defined as a source language for the final project in courses on compiler construction, as most students are already
familiar with the Python language by the time they take such a course. By restricting the language to a statically typed fragment, most
student compilers can produce code that usually even outperforms the CPython reference interpreter. In a similar manner, Restricted Python 
(RPython)~\cite{ancona_rpython_2007} was defined as a strictly typed subset of Python, which removes some of the more dynamic features of
the language. As part of the PyPy project~\cite{rigo_pypys_2006}, this allows RPython to be more easily compiled to common language 
runtimes, such as JVM and CLI/.NET.

We are aware of only one other work~\cite{sunderraman_functional_2025} specifically targetting a functional subset of Python, with
the express intent of using it as a teaching tool. While we share a similar goal, our language fragment encloses more of the overall
Python language, as their work doesn't allow for control structures such as conditionals or loops.

While it is not our goal to change Python towards a more functional programming language, it is interesting to notice that there
are movements within the Python community for inclusion of programming principles commonly found in functional languages. An example
of this is the inclusion of \emph{lambda expressions} already in Python 1.0, and a more recent PEP draft~\cite{johnson_pep_2025}
advocating for \emph{deeply immutable objects}. 
\section{Conclusions and future work}
\label{sec:conclusion}

This section summarises the contributions and discusses future work.


\bibliographystyle{ACM-Reference-Format}
\bibliography{tex/refs,tex/additional-refs}

\end{document}
